% Canzoniere di Clan - template. I segnaposto tra doppi angoli (maiuscoli)
% sono sostituiti da scripts/generate_latex.py. Compilare con pdflatex (due passate) o latexmk.
\documentclass[a5paper,landscape,10pt,twoside]{article}

\usepackage[utf8]{inputenc}
\usepackage[T1]{fontenc}
\usepackage[italian]{babel}
\usepackage{lmodern}
\usepackage[a5paper,landscape,twoside,
  left=15mm,right=15mm,top=12mm,bottom=12mm,bindingoffset=15mm,
  includefoot,nohead,footskip=6mm]{geometry}
\usepackage{multicol}
\usepackage{graphicx}
\usepackage{xcolor}
\usepackage{url}
\usepackage{microtype}

\setlength{\parindent}{0pt}
\setlength{\parskip}{0pt}
\setlength{\columnsep}{8mm}
\raggedbottom
\sloppy
\hyphenpenalty=1000

\definecolor{accordo}{gray}{0.25}
\definecolor{brLC}{HTML}{F2C230}
\definecolor{brEG}{HTML}{3F9A47}
\definecolor{brRS}{HTML}{C0392B}
\definecolor{fiamma}{HTML}{D9822B}

% Interlinea ampia: lascia spazio agli accordi sopra il testo.
\makeatletter
\renewcommand{\normalsize}{\@setfontsize\normalsize{10}{17}}

% Stile pagina: solo numero in basso, sul lato esterno.
\def\ps@canz{%
  \def\@oddhead{}\def\@evenhead{}%
  \def\@oddfoot{\hfill\small\thepage}%
  \def\@evenfoot{\small\thepage\hfill}}

% \cantotesta{id}{titolo}{categoria}{autore}: nuova pagina + etichetta per gli indici
% \fiamme{n}: intensita' 1-5 come cinque barrette piene/vuote
\newcommand{\barra}{\textcolor{fiamma}{\rule{1.1mm}{2.6mm}}\,}
\newcommand{\barraoff}{\textcolor{black!20}{\rule{1.1mm}{2.6mm}}\,}
\newcommand{\fiamme}[1]{{\@tempcnta=\z@
  \@whilenum\@tempcnta<5\do{\advance\@tempcnta\@ne
    \ifnum\@tempcnta>#1\barraoff\else\barra\fi}}}
% Tag di branca: solo un quadratino colorato (giallo LC, verde EG, rosso RS)
\newcommand{\tagbranca}[1]{\textcolor{#1}{\rule{2.4mm}{2.4mm}}\,}
\newcommand{\bLC}{\tagbranca{brLC}}
\newcommand{\bEG}{\tagbranca{brEG}}
\newcommand{\bRS}{\tagbranca{brRS}}

% \cantotesta{id}{titolo}{categoria}{autore}{tag branche}{intensita'}
\newcommand{\cantotesta}[6]{%
  \clearpage
  \def\@currentlabel{#1}\label{canto:#1}%
  \noindent\hangindent=14mm\hangafter=1
  \hbox to 14mm{\fontsize{22}{24}\selectfont\bfseries #1\hss}%
  {\fontsize{15}{18}\selectfont\bfseries #2\par}
  \noindent{\small\scshape #3\ifx\relax#4\relax\else\ \textperiodcentered\ \textit{#4}\fi}\hfill{#5\ #6}\par
  \vspace{1mm}\hrule height 0.4pt\vspace{3mm}}
\makeatother
\pagestyle{canz}

% Riga di testo con accordi: l'accordo sta sopra la sillaba che lo segue.
\newcommand{\ch}[2]{\leavevmode\vbox{\offinterlineskip
  \hbox{\fontsize{8}{9}\selectfont\bfseries\color{accordo}#1\kern.3em}%
  \kern.45ex
  \hbox{#2}}}
\newcommand{\rigo}[1]{{\hangindent=3mm\hangafter=1\noindent #1\par}}
\newcommand{\strofa}{\par\vspace{2.4mm}}
\newcommand{\commento}[1]{{\par\small\itshape\color{gray}#1\par\vspace{1mm}}}
\newenvironment{ritornello}{\par\begingroup\itshape\leftskip=3mm\relax}{\par\endgroup}

% Indici
\newcommand{\titolosezione}[1]{\clearpage
  {\noindent\fontsize{16}{18}\selectfont\bfseries #1\par}
  \vspace{1mm}\hrule height 0.4pt\vspace{3mm}}
\newcommand{\lettera}[1]{\par\vspace{1.5mm}{\bfseries\color{accordo}#1}\par\nopagebreak}
\newcommand{\categoriaindice}[1]{\par\vspace{2mm}{\scshape\bfseries #1}\par\nopagebreak}
\newcommand{\voce}[2]{\noindent\hangindent=2.4em\hangafter=1
  \hbox to 2em{\hfil\bfseries #1}\hskip .4em #2\nobreak
  \leaders\hbox to .7em{\hss.\hss}\hfill\nobreak
  \hbox to 2em{\hfill\pageref{canto:#1}}\par}

\newcommand{\sottocategoriaindice}[1]{\par\vspace{.8mm}{\small\itshape #1}\par\nopagebreak}

\begin{document}

% ---- Frontespizio (numerazione romana, nessun numero stampato)
\pagenumbering{roman}
\thispagestyle{empty}
\begin{center}
  \vspace*{\fill}
  {\fontsize{40}{44}\selectfont\bfseries Canzoniere\par}
  \vspace{4mm}
  {\Large Clan - Gruppo Scout\par}
  \vspace{10mm}
  {\large\itshape Un regalo di Partenza, per i miei compagni di Clan.\par}
  \vspace*{\fill}
  {\small 2026\par}
\end{center}
% Pagina di dedica. BOZZA da personalizzare: sostituisci il testo e la firma.
\clearpage
\thispagestyle{empty}
\vspace*{\fill}
\begin{center}
  \begin{minipage}{0.8\textwidth}
    \centering\itshape\large
    Ai miei compagni di Clan.\par
    \vspace{5mm}
    \normalsize
    Ho camminato con voi per tanta strada, e quasi ogni passo l'abbiamo cantato.
    Questo libro raccoglie le canzoni che ci hanno tenuti insieme: quelle gridate
    in marcia, quelle sussurrate attorno al fuoco, quelle che non abbiamo mai
    finito di imparare.\par
    \vspace{3mm}
    Portatelo con voi, lasciate che si consumi, e quando io non ci sar\`o
    cantatene una anche per me.\par
    \vspace{6mm}
    Buona strada.\par
    \vspace{3mm}
    \upshape [Nome e totem]
  \end{minipage}
\end{center}
\vspace*{\fill}


\titolosezione{Come usare questo canzoniere}
\begingroup\small
\begin{multicols}{2}
\textbf{Struttura.} I canti sono numerati e raggruppati in sezioni, nell'ordine:
Canti associativi, Canti cerimoniali, Canti di lupettismo, Fuoco di bivacco, Canti alpini, Canti liturgici, Canti del risveglio. Ogni canto ha la sua pagina: il numero grande accanto al titolo
è quello da cercare, sia qui sia nel sito. In fondo trovi l'indice alfabetico
e l'indice per sezione.

\medskip
\textbf{Accordi.} Sono scritti sopra la sillaba in cui cambiano. Il testo in
\textit{corsivo} è il ritornello.

\medskip
\textbf{Online.} Sull'ultima pagina c'è un QR code: porta al sito con l'audio di
ogni canto, gli spartiti in PDF e la modalità notte per il bivacco.

\columnbreak
\textbf{Branche.} A destra del titolo, le branche per cui il canto nasce o è
più adatto:

\smallskip
\bLC\ lupetti e coccinelle \quad \bEG\ esploratori e guide \quad \bRS\ rover e scolte

\medskip
\textbf{Intensità.} Quanta energia chiede il canto, da 1 a 5:

\smallskip
\fiamme{1}\ raccolto, quasi sussurrato\\
\fiamme{3}\ corale, a voce piena\\
\fiamme{5}\ esplosivo, con gesti e balli

\medskip
Un canto senza tag è per tutti.
\end{multicols}
\endgroup


\titolosezione{Il fuoco di bivacco}
\begingroup\small
Un buon fuoco ha un andamento a onda: sale, tocca il punto più alto e poi scende fino al silenzio.
Scegliere i canti seguendo l'intensità (da \fiamme{1} a \fiamme{5}) aiuta a non bruciare tutto subito.

\medskip
\begin{tabular}{@{}l@{\hspace{4mm}}p{0.72\textwidth}@{}}
\textbf{Apertura} & \fiamme{3}\ Si accende il fuoco, ci si riunisce. Canti noti, facili da attaccare, che fanno sentire tutti dentro.\\[1.5mm]
\textbf{Salita} & \fiamme{4}\ Bans, canoni e canti con gesti: si scioglie il gruppo e si prende ritmo.\\[1.5mm]
\textbf{Picco} & \fiamme{5}\ Il momento più acceso, breve. Canti gridati, balli, voci che si sovrappongono.\\[1.5mm]
\textbf{Discesa} & \fiamme{3}\ Si abbassa piano la voce e il ritmo; si torna a cantare insieme, non uno contro l'altro.\\[1.5mm]
\textbf{Raccoglimento} & \fiamme{2}\ Canti lenti, di strada e di ricordo. Il fuoco è basso, si canta quasi sottovoce.\\[1.5mm]
\textbf{Chiusura} & \fiamme{1}\ Un canto semplice o una preghiera, poi silenzio. Non si applaude: si lascia che resti.
\end{tabular}

\medskip
\textit{Regola d'oro:} chi guida legge il fuoco, non la scaletta. Se il gruppo è stanco, si scende prima.
\endgroup


\titolosezione{Cantare insieme}
\begingroup\small
\begin{multicols}{2}
\textbf{1. Il tono.} Chi attacca sceglie una nota che stia comoda a tutti: meglio un po' bassa che
troppo alta. Se la prima strofa è faticosa, si riparte da capo, più in basso.

\smallskip
\textbf{2. Un solo attacco.} Si dà il tempo con un respiro o con la mano e si parte insieme.
Cantare a ritmi diversi è il modo più rapido per spegnere un canto.

\smallskip
\textbf{3. Ascoltare.} Si canta in modo da sentire il vicino. Se non lo senti più, stai coprendo
la sua voce.

\smallskip
\textbf{4. Il respiro.} Si respira alla fine delle frasi, non a metà parola. Chi respira in
momenti diversi tiene su il canto senza buchi.

\columnbreak
\textbf{5. Canoni e voci.} Nei canoni ogni gruppo tiene la propria parte e non segue quello
accanto. Si comincia a due voci, poi si aggiungono le altre.

\smallskip
\textbf{6. Bans.} Si insegnano un pezzo alla volta: prima i gesti, poi le parole, poi tutto
insieme. Chi guida sta davanti e li fa per primo.

\smallskip
\textbf{7. Chitarra.} Sta sotto le voci, non sopra. Un giro semplice e regolare vale più di uno
elaborato; con il capotasto si abbassa il tono senza cambiare accordi.

\smallskip
\textbf{8. Senza strumenti.} Un canto deve reggere anche a cappella: se si perde l'intonazione,
ci si ferma, si ritrova la nota e si riparte.
\end{multicols}
\endgroup


% ---- Canti (la numerazione araba riparte da 1: pagina stampata = campo "pagina")
\pagenumbering{arabic}
\cantotesta{1}{Inno d'Italia}{Canti associativi}{G. Mameli - M. Novaro}{}{}
\begin{multicols}{2}
\rigo{\ch{C}{Fratelli} d'Italia,}
\rigo{l'Italia s'è \ch{G}{desta,}}
\rigo{dell'elmo di \ch{C}{Scipio}}
\rigo{s'è cinta la \ch{G}{testa.}}
\rigo{Dov'è la Vit\ch{C}{toria?}}
\rigo{Le porga la \ch{G}{chioma,}}
\rigo{ché schiava di \ch{C}{Roma}}
\rigo{Iddio la cre\ch{G}{ò.}}
\strofa
\begin{ritornello}
\rigo{\ch{C}{Stringiamci} a coorte,}
\rigo{siam \ch{G}{pronti} alla morte,}
\rigo{siam \ch{C}{pronti} alla morte,}
\rigo{l'Italia chia\ch{G}{mò.} \ch{C}{Sì!}}
\end{ritornello}
\end{multicols}

\cantotesta{2}{Addio lupetto}{Canti cerimoniali}{}{\bLC}{}
\begin{multicols}{2}
\rigo{Formiamo un grande cerchio}
\rigo{con le mani nelle man;}
\rigo{fratelli c'è un lupetto}
\rigo{che oggi il branco lascerà.}
\strofa
\rigo{Rit. Ma noi vogliamo dirgli ancor}
\rigo{che un dì ci rivedrem.}
\rigo{Arrivederci, fratellin,}
\rigo{presto ti seguirem.}
\strofa
\rigo{Tu pure come Mowgli}
\rigo{lasci il branco e vai lontan,}
\rigo{tu pure come Mowgli}
\rigo{sarai uomo nel doman.}
\strofa
\rigo{Ricordati del branco}
\rigo{noi di te fieri sarem}
\rigo{nessun ti scorderà,}
\rigo{fratello e un dì ti seguirem.}
\rigo{> Indice}
\end{multicols}

\cantotesta{3}{Canto della promessa del lupetto}{Canti cerimoniali}{}{\bLC}{}
\begin{multicols}{2}
\rigo{Davanti a Te m'impegno}
\rigo{con tutto il cuor}
\rigo{diventerò tuo paggio}
\rigo{o mio Signor.}
\strofa
\rigo{Io voglio amarti sempre}
\rigo{sempre più}
\rigo{e la promessa mia}
\rigo{proteggi Tu.}
\strofa
\rigo{Voglio esser Tuo lupetto}
\rigo{per divenir}
\rigo{più tardi esploratore}
\rigo{per servir.}
\strofa
\rigo{Io voglio amarti sempre}
\rigo{sempre più}
\rigo{e la promessa mia}
\rigo{proteggi Tu.}
\strofa
\rigo{A Te Gesù che amo}
\rigo{obbedirò}
\rigo{a me dimenticando}
\rigo{ti ascolterò}
\strofa
\rigo{Io voglio amarti sempre}
\rigo{sempre di più}
\rigo{e la promessa mia}
\rigo{proteggi Tu.}
\strofa
\rigo{> Indice}
\end{multicols}

\cantotesta{4}{Canto della promessa}{Canti cerimoniali}{J. Sevin - tradizionale}{\bEG}{}
\begin{multicols}{2}
\rigo{\ch{Re}{Dinanzi} a voi m'impegno sul mio onor   \ch{La}{}}
\rigo{e vog\ch{Re}{lio} esserne degno per te o \ch{La}{Signor}         \ch{Re}{}}
\strofa
\rigo{Rit.}
\rigo{SOL             RE         LA7             RE         7}
\rigo{La giusta e retta via, mostrami Tu}
\rigo{SOL            RE     LA    7             RE}
\rigo{e la promessa mia accogli o Gesù.}
\strofa
\rigo{Fedele al tuo volere sempre sarò,}
\rigo{di Patria il mio dovere adempirò.}
\strofa
\rigo{Rit.}
\strofa
\rigo{Apostolo tuo sono per tuo amor,}
\rigo{agli altri di me dono vo’ fare ognor.}
\strofa
\rigo{Rit.}
\strofa
\rigo{Leale alla mia Legge sempre sarò,}
\rigo{se la tua man mi regge io manterrò!}
\strofa
\rigo{Rit.}
\strofa
\rigo{> Indice}
\end{multicols}

\cantotesta{5}{Canto delle Scolte d'Assisi}{Canti cerimoniali}{}{\bRS}{}
\begin{multicols}{2}
\rigo{Squilla!}
\rigo{Squilla la tromba che già il giorno finì}
\rigo{già del coprifuoco la canzone salì}
\rigo{Su, scolte, alle torri guardie armate ola'}
\rigo{attente in silenzio vigilate.}
\strofa
\rigo{O nostri santi che in cielo esultate}
\rigo{Vergini sante gloriose e beate}
\rigo{noi v'invochiamo.}
\strofa
\rigo{Questa città con il vostro amore}
\rigo{guardate contro il nemico che l'anima}
\rigo{tiene contro la morte che subito viene}
\rigo{in ogni cuor sia pace e bene}
\rigo{sia tregua ad ogni dolor. Pace.}
\strofa
\rigo{Squilla squilla la tromba che già il giorno finì}
\rigo{già del coprifuoco la canzone salì.}
\strofa
\rigo{Su, scolte, alle torri guardie armate ola'}
\rigo{attente in silenzio vigilate.}
\strofa
\rigo{Attente, scolte, vigilate.}
\strofa
\rigo{> Indice}
\end{multicols}

\cantotesta{6}{Servire}{Canti cerimoniali}{}{\bRS}{}
\begin{multicols}{2}
\rigo{Rit: Servire è parola che parla col cuore}
\rigo{servire è parola lucente d'affetto}
\rigo{servire è un concetto fraterno d'amore}
\rigo{che unisce, stupisce, che parla da se}
\strofa
\rigo{La strada fraterna è cosparsa d'amore}
\rigo{l'amore è paziente, è benigno l'amore}
\rigo{se dessi tutti i miei averi di cuore}
\rigo{e non dessi me stesso a che serve l'amore}
\strofa
\rigo{Rit.:}
\strofa
\rigo{L'amore non invidia, non riesce a vantare}
\rigo{non gonfia, non manca rispetto di dare}
\rigo{non cerca interesse, l'amore non si adira}
\rigo{ma cerca ed aspira soltanto a servire.}
\strofa
\rigo{Rit.:}
\strofa
\rigo{Se ami la feste, il bosco, il ruscello}
\rigo{se cerchi l'amore che cerca un fanciullo}
\rigo{sel male, sai, il conto non devi tenere}
\rigo{ma devi soltanto pensare a servire.}
\strofa
\rigo{Rit.:}
\strofa
\rigo{Fratello, fratello, servire è parola}
\rigo{pronunciata una volta è ricordata per sempre}
\rigo{sui campi, nei boschi da rover e da anziani}
\rigo{in mezzo alla folla da oggi al domani.}
\strofa
\rigo{> Indice}
\end{multicols}

\cantotesta{7}{A caccia andiamo}{Canti di lupettismo}{}{\bLC}{}
\begin{multicols}{2}
\rigo{O lupi a caccia andiamo,}
\rigo{corriamo e pronti siam}
\rigo{sorriso lieto abbiamo}
\rigo{"del nostro meglio" urliam}
\rigo{e meglio urliam.}
\strofa
\rigo{Sì, bisognerà, bisognerà che vada.}
\rigo{Sì, bisognerà, bisognerà partir.}
\rigo{Sì, bisognerà, bisognerà che vada.}
\rigo{Sì, bisognerà, bisognerà partir.}
\strofa
\rigo{In mezzo alla natura}
\rigo{fratelli noi sarem.}
\rigo{Lupetti all'aria pura}
\rigo{felici correrem e correrem.}
\strofa
\rigo{Sì, bisognerà...}
\strofa
\rigo{Rientrati nella tana}
\rigo{ricorderemo, sì,}
\rigo{la caccia ormai lontana,}
\rigo{cantando ognor così, ognor così.}
\strofa
\rigo{Sì, bisognerà...}
\strofa
\rigo{> Indice}
\end{multicols}

\cantotesta{8}{Attorno alla rupe}{Canti di lupettismo}{}{\bLC}{}
\begin{multicols}{2}
\rigo{Attorno alla rupe orsù lupi andiam}
\rigo{d'Akela e Baloo or le voci ascoltiam}
\rigo{del branco la forza in ciascun lupo sta}
\rigo{del lupo la forza del branco sarà.}
\strofa
\rigo{Rit.: uulla ullalala uullalala}
\rigo{del lupo la forza nel branco sarà (2v.)}
\strofa
\rigo{Or quando si sente un richiamo lontan}
\rigo{risponde un vicino più forte ulular}
\rigo{a quel che la legge fedel seguirà}
\rigo{a lui buona caccia si ripeterà}
\strofa
\rigo{Rit.: uulla ullalala uullalala}
\rigo{a lui buona caccia si ripeterà}
\strofa
\rigo{Il debole cucciolo via via crescerà}
\rigo{e forte ed ardito col clan caccerà}
\rigo{finché sulla rupe un dì porterà}
\rigo{la pelle striata del vile Shere Khan.}
\strofa
\rigo{Rit.: uulla ullalala uullalala}
\rigo{la pelle striata del vile Shere Khan (2v.)}
\strofa
\rigo{> Indice}
\end{multicols}

\cantotesta{9}{Branco di Akela}{Canti di lupettismo}{}{\bLC}{}
\begin{multicols}{2}
\rigo{Akela aveva un branco felice e fier,}
\rigo{sestiglie rossi e fulvi, pezzati e ner,}
\rigo{furbi i rossi in verità, con la loro abilità}
\rigo{nel seguir la traccia di Akela!.}
\strofa
\rigo{Ridenti anche i pezzati felici ognor}
\rigo{perché sestiglia in gamba è quella lor.}
\rigo{Ecco adesso venir giù dalla caccia di Baloo}
\rigo{vincitori i neri pien di virtù.}
\strofa
\rigo{Audaci e agili sono i fulvi ancor}
\rigo{i primi della corsa e allegri in cor.}
\rigo{La miglior quale sarà che per prima arriverà}
\rigo{a conquistare la preda di Akela.}
\strofa
\rigo{> Indice}
\end{multicols}

\cantotesta{10}{Buona caccia}{Canti di lupettismo}{}{\bLC}{}
\begin{multicols}{2}
\rigo{Partiamo col branco in caccia buona caccia, buona caccia;}
\rigo{partiamo col branco in caccia buona caccia fratellin.}
\strofa
\rigo{Cuor leale, lingua cortese fanno strada nella giungla;}
\rigo{cuor leale, lingua cortese nella giunglastrada fan.}
\strofa
\rigo{Siamo d'uno stesso sangue fratellino tu ed io;}
\rigo{siamo d'uno stesso sangue fratellino io e te.}
\strofa
\rigo{E' le legge della giungla vecchia e vera come il cielo}
\rigo{senza legge non c'e' vero non c'e' lupo che vivra'.}
\strofa
\rigo{Grande e' il bosco e tenebroso ed il lupo e' piccolino;}
\rigo{Baloo, l'orso coraggioso, viene a caccia insieme a te.}
\strofa
\rigo{> Indice}
\end{multicols}

\cantotesta{11}{Caccia al sambhur}{Canti di lupettismo}{}{\bLC}{}
\begin{multicols}{2}
\rigo{Nella valletta chiara s'aggirano i lupi}
\rigo{annusano le tracce del sambhur veloce}
\rigo{seguono Akela a caccia tra aspri dirupi}
\rigo{e da lontano l'eco ripete la voce.}
\strofa
\rigo{E salta, lupo, salta, salta}
\rigo{e salta sulla preda}
\rigo{Shere Khan che non ti veda}
\rigo{che te la vuol levar! (bis)}
\strofa
\rigo{> Indice}
\end{multicols}

\cantotesta{12}{Canto del risveglio}{Canti di lupettismo}{}{\bLC}{}
\begin{multicols}{2}
\rigo{MI b                      SI b             MI b}
\rigo{Lupo salta su che già l'ora è suonata !}
\rigo{SI b     MI b}
\rigo{Lupo salta su, cominciato è il dì.}
\rigo{DO- SI b         MI b}
\rigo{Tra poco al ruscelletto fra le rocce correremo}
\rigo{DO m          SI b         MI b}
\rigo{e lieti all'acqua fresca tutti insiem ci laveremo.}
\strofa
\rigo{La giungla già ci attende}
\rigo{presto in caccia partiremo,}
\rigo{la pista in mezzo al bosco}
\rigo{pronti e svelti seguiremo.}
\strofa
\rigo{> Indice}
\end{multicols}

\cantotesta{13}{Canto della sera}{Canti di lupettismo}{}{\bLC}{}
\begin{multicols}{2}
\rigo{Or scende la sera e noi siamo quaggiù}
\rigo{tutto il Branco è in preghiera}
\rigo{e Ti invoca, o Gesù.}
\rigo{Dall'alto del cielo proteggici Tu}
\rigo{ed il nostro branco benedici, o Gesù.}
\rigo{Quel che oggi abbiam fatto e le nostre B.A.}
\rigo{ecco a Te presentiamo con tutta lealtà.}
\strofa
\rigo{Proteggici a casa la mamma e il papà}
\rigo{fa' che noi meritiamo la loro bontà.}
\rigo{Silente è la giungla nel giorno che muor}
\rigo{su tutti discenda la gioia e l'amor.}
\rigo{Dall'alto del cielo proteggici Tu}
\rigo{ed il nostro branco benedici o Gesù.}
\strofa
\rigo{> Indice}
\end{multicols}

\cantotesta{14}{Danza dei cani rossi}{Canti di lupettismo}{}{\bLC}{}
\begin{multicols}{2}
\rigo{Rah, rah, rah! il cane rosso è qua.}
\rigo{Arriva giù, non lascia più la traccia di Wontolla.}
\strofa
\rigo{Rah, rah, rah! Il cane rosso è qua.}
\rigo{La rabbia sì lo inferocì, Wontolla lo annunciò.}
\strofa
\rigo{Or Bagheera, Kaa e Baloo alla lotta pronti son;}
\rigo{al comando di Akela coi suoi lupi Mowgli vien.}
\strofa
\rigo{Rah, rah, rah! Il cane rosso è qua.}
\rigo{Di peso fu portato giù dall'acqua del Waingunga.}
\strofa
\rigo{Rah, rah, rah! Il cane rosso è qua.}
\rigo{E quando là approderà il branco pronto è già.}
\strofa
\rigo{Tutti i lupi di Seeonee stan lottando con ardor}
\rigo{ed il Dhole invasor alla fin sconfiggeran.}
\strofa
\rigo{> Indice}
\end{multicols}

\cantotesta{15}{Danza del fiore rosso}{Canti di lupettismo}{}{\bLC}{}
\begin{multicols}{2}
\rigo{C'era un tempo un tigre perfida e sleal}
\rigo{nella giungla faceva tanto tanto mal}
\rigo{ma poi Mowgli il fiore rosso}
\rigo{alla rupe ci portò}
\rigo{ed il pelo ai suoi nemici}
\rigo{con un ramo bruciacchiò.}
\rigo{Fiore rosso avvampa ancor}
\rigo{tutto il branco ne ha terror.}
\strofa
\rigo{Mowgli prese pel mento}
\rigo{il muso di Shere-Khan}
\rigo{mentre il fior sbocciava}
\rigo{nella salda man;}
\rigo{tigre zoppa della giungla}
\rigo{la tua pelle brucerò}
\rigo{sulla rupe del Consiglio}
\rigo{se non fuggi stenderò.}
\rigo{Fiore rosso avvampa ancor}
\rigo{anche Shere-Khan ne ha terror.}
\strofa
\rigo{> Indice}
\end{multicols}

\cantotesta{16}{Danza della tregua dell'acqua}{Canti di lupettismo}{}{\bLC}{}
\begin{multicols}{2}
\rigo{Lupo la parola d'Hathi ascolterà}
\rigo{quando nella giungla tregua annunzierà.}
\strofa
\rigo{Rit. Tra la la la la ra la la la}
\rigo{tra la la la la ra la la la}
\rigo{tra la la la la ra la la la}
\rigo{tra la la la la ra la la la.}
\strofa
\rigo{Tigre col cerbiatto all'acqua insieme andrà}
\rigo{se nella Waingunga la roccia si vedrà.}
\strofa
\rigo{Acqua fresca e viva Baloo berrà con te...}
\rigo{ma sulla sorgente una macchia scura c'é.}
\strofa
\rigo{Gridano d'orrore Bagheera, Kaa, Baloo:}
\rigo{via da questa fonte, non ci tornare più!!!.}
\strofa
\rigo{E Shere-khan che arriva... d'umano sangue ha}
\rigo{zanne rosse orrende, ed ora a bere va.}
\strofa
\rigo{> Indice}
\end{multicols}

\cantotesta{17}{Fratelli alla candida luna}{Canti di lupettismo}{}{\bLC}{}
\begin{multicols}{2}
\rigo{RIT: Fratelli alla candida luna cantiamo la nostra canzone più bella:}
\rigo{tra faggi ed abeti, che gioia veniamo la nostra canzone a cantar.}
\strofa
\rigo{1- "Una via sola è vera" ci disse il lupo anziano}
\rigo{e dietro a lui noi siamo la traccia a seguitar.}
\rigo{Costanti alla ricerca, gioiosi nel cammino,}
\rigo{erriam fino al mattino seguendo il lupo anzian.}
\strofa
\rigo{2- La via è troppo scoscesa la traccia già si perde.}
\rigo{Solo non si disperde chi segue il lupo anzian.}
\rigo{Lupetto non lasciare oh, non lasciare il branco}
\rigo{il pelo è ancora bianco seguiamo il lupo anzian.}
\strofa
\rigo{> Indice}
\end{multicols}

\cantotesta{18}{Il canto del cuculo}{Canti di lupettismo}{}{\bLC}{}
\begin{multicols}{2}
\rigo{Sentiam nella foresta il cuculo cantar}
\rigo{ai piedi d'una quercia lo stiamo ad ascoltar.}
\strofa
\rigo{Cucu cucu...}
\strofa
\rigo{La notte è tenebrosa non c'è chiaror lunar}
\rigo{sentiam nel fitto bosco i lupi ulular.}
\strofa
\rigo{Ahu, ahu...}
\strofa
\rigo{> Indice}
\end{multicols}

\cantotesta{19}{Il lupo ferito}{Canti di lupettismo}{}{\bLC}{}
\begin{multicols}{2}
\rigo{Bagheera che vide narrò}
\rigo{d'un lupo dal giovane ardir}
\rigo{che il branco alla caccia lasciò}
\rigo{e volle l'istinto seguir.}
\strofa
\rigo{Rit. Lupo lupo torna nel branco}
\rigo{a cacciar con noi vieni vieni}
\rigo{vieni il grand'urlo a lanciar.}
\strofa
\rigo{La legge che insegna il dover}
\rigo{la legge che Akela ululò}
\rigo{la legge che un branco fa fier}
\rigo{il lupo randagio scordò.}
\strofa
\rigo{Ma quando dall'ombra sbucò}
\rigo{il toro selvaggio e guerrier}
\rigo{invano quel lupo saltò}
\rigo{ferito ululò sul sentier.}
\strofa
\rigo{E il lupo l'errore capì}
\rigo{tornò più veloce nel pian}
\rigo{e corse alla rupe e finì}
\rigo{fra le zampe di un lupo anzian.}
\strofa
\rigo{> Indice}
\end{multicols}

\cantotesta{20}{Il risveglio della giungla}{Canti di lupettismo}{}{\bLC}{}
\begin{multicols}{2}
\rigo{La luna è bianca e scuro è l'orizzonte}
\rigo{già tutto tace, sol canta la fonte.}
\strofa
\rigo{Lassù nel cielo bruno tutto solo}
\rigo{già Mang, il pipistrello, intreccia il volo.}
\strofa
\rigo{La giungla si ridesta piano piano}
\rigo{ed echeggia di lontano già il richiamo:}
\rigo{Buona caccia. Buona caccia.}
\strofa
\rigo{a tutti quelli che rispettano la Legge;}
\rigo{Buona caccia Buona caccia Buona caccia!}
\strofa
\rigo{> Indice}
\end{multicols}

\cantotesta{21}{L'ultimo canto di Akela}{Canti di lupettismo}{}{\bLC}{}
\begin{multicols}{2}
\rigo{Il lupo Wontolla non caccerà più}
\rigo{Akela da un pezzo è disteso laggiù}
\rigo{e s'ode d'intorno stormire di fronde}
\rigo{sommessa sussurra la giungla un sospir.}
\strofa
\rigo{Rit. E muoio contento fratello a sentir}
\rigo{che sei qui vicino a vedermi morir.}
\strofa
\rigo{E Mowgli il ranocchio lo venne a cercar}
\rigo{tra i lupi caduti l'udì rantolar}
\rigo{alzata la testa e le spalle di lì}
\rigo{a Mowgli in ginocchio diceva così.}
\strofa
\rigo{Già l'ultimo colpo il mio dente vibrò}
\rigo{tra poco la giungla lasciare dovrò.}
\rigo{Ma muoio contento, fratello, a sentir}
\rigo{che sei qui vicino a vedermi morir!}
\strofa
\rigo{> Indice}
\end{multicols}

\cantotesta{22}{Lupo salta su}{Canti di lupettismo}{}{\bLC}{}
\begin{multicols}{2}
\rigo{MI b                          SI b          MI b}
\rigo{Lupo salta su che già l’ora è suonata !}
\rigo{SI b     MI b}
\rigo{Lupo salta su, cominciato è il dì.}
\rigo{DO- SI b        MI b}
\rigo{Tra poco al ruscelletto fra le rocce correremo}
\rigo{DO m        SI b         MI b}
\rigo{e lieti all’acqua fresca tutti insiem ci laveremo.}
\strofa
\rigo{La giungla già ci attende}
\rigo{presto in caccia partiremo,}
\rigo{la pista in mezzo al bosco}
\rigo{pronti e svelti seguiremo.}
\strofa
\rigo{> Indice}
\end{multicols}

\cantotesta{23}{Salta o lupo}{Canti di lupettismo}{}{\bLC}{}
\begin{multicols}{2}
\rigo{Il sentiero di caccia è aperto innanzi a te}
\rigo{ascolta l'urlo, s'ode da lontan.}
\rigo{Il favor della giungla ti accompagnerà}
\rigo{sarà il tuo passo forte ardito e fier.}
\strofa
\rigo{Rit. Salta o lupo se tu vuoi ogni preda}
\rigo{sarà presto ai piedi tuoi}
\rigo{giuste tracce seguirai}
\rigo{e i segreti della giungla scoprirai.}
\strofa
\rigo{Se un buon lupo presto tu vuoi diventar}
\rigo{le Massime e la Legge non scordare,}
\rigo{la parola di Akela sia vicino a te}
\rigo{osserva sempre i detti di Baloo.}
\strofa
\rigo{Occhio aperto alle stelle dritto sa guardar}
\rigo{Bagheera a caccia ti può consigliar.}
\rigo{Del tuo branco la forza in te stesso sta}
\rigo{e la sua forza il branco a te darà.}
\strofa
\rigo{Quando il branco alla rupe si radunerà}
\rigo{di nuove imprese ancor potrem parlar.}
\rigo{Quella Legge che uniti sempre manterrà}
\rigo{il nostro vanto e il tuo onor sarà.}
\strofa
\rigo{> Indice}
\end{multicols}

\cantotesta{24}{Siam nella vasta giungla}{Canti di lupettismo}{}{\bLC}{}
\begin{multicols}{2}
\rigo{Siam nella vasta giungla}
\rigo{siam sulla pista di Akela}
\rigo{lieti cantando, sempre cantando}
\rigo{Lupi siam.}
\strofa
\rigo{> Indice}
\end{multicols}

\cantotesta{25}{Ula, ula, ula}{Canti di lupettismo}{}{\bLC}{}
\begin{multicols}{2}
\rigo{Ula ula ula}
\rigo{è sera e stanchi siam}
\rigo{ula ula ula}
\rigo{domani lieti sarem}
\rigo{ula ula ula}
\rigo{tra poco dormirem}
\rigo{per tutta la notte}
\rigo{noi ci riposerem}
\rigo{ula ula ula}
\rigo{il branco dorme già}
\rigo{ula ula ula}
\rigo{la giungla tacerà}
\rigo{ula ula ula}
\rigo{Gesù benedirà}
\rigo{questa nostra tana}
\rigo{e ci sorriderà...}
\strofa
\rigo{> Indice}
\end{multicols}

\cantotesta{26}{Un bravo lupo}{Canti di lupettismo}{}{\bLC}{}
\begin{multicols}{2}
\rigo{Un bravo lupo io voglio diventar}
\rigo{e la promessa per sempre rispettar}
\rigo{gentile e più cortese con tutti io sarò}
\rigo{la Buona Azion sempre farò.}
\strofa
\rigo{Akela oh Akela oh, Te lo prometto,}
\rigo{più in gamba io sarò.}
\strofa
\rigo{Caro Baloo, caro Baloo}
\rigo{io la legge osserverò di più. ( 2 volte )}
\strofa
\rigo{Bagheera oh, Bagheera oh}
\rigo{te lo prometto, in caccia io verrò. ( 2 volte )}
\strofa
\rigo{Mio caro Kaa, mio caro Kaa}
\rigo{te lo prometto, farò tante B.A. ( 2 volte )}
\strofa
\rigo{> Indice}
\end{multicols}

\cantotesta{27}{Zampe agili e silenti}{Canti di lupettismo}{}{\bLC}{}
\begin{multicols}{2}
\rigo{Zampe agili e silenti}
\rigo{zanne aguzze e rilucenti}
\rigo{acuto l'occhio scorge nell'oscurità}
\rigo{orecchie tese odono il vento che va.}
\rigo{D'un lupo sono i segni}
\rigo{d'un forte cacciator.}
\strofa
\rigo{> Indice}
\end{multicols}

\cantotesta{28}{La leggenda del fuoco}{Fuoco di bivacco · Apertura}{}{\bEG}{}
\begin{multicols}{2}
\rigo{\ch{Do}{Splende} il fuoco nel cerchio dell’esplorato\ch{Sol}{r}   \ch{Do}{}}
\rigo{ascoltate la voce della f\ch{Sol}{iamma} d’or. \ch{Do}{}}
\strofa
\rigo{RIT.}
\rigo{Sali al c\ch{Sol}{iel} fiamma leggera  \ch{Do}{}}
\rigo{del gran fuoc\ch{Fa}{o} caldo e buon     \ch{Sol}{}}
\rigo{sotto i pini alla brughie\ch{Do}{ra}}
\rigo{sali in a\ch{Fa}{lto} sali ancor \ch{Sol}{}}
\rigo{sali in alto sali anc\ch{Do}{or,}}
\rigo{fuoco dell’es\ch{Fa}{plorator.}  \ch{Sol}{}}
\strofa
\rigo{Ero un principe un giorno perfido e sleal}
\rigo{e spargevo d’intorno il terrore e il mal.}
\strofa
\rigo{RIT.}
\strofa
\rigo{D’un gran mago l’incanto tosto mi punì}
\rigo{e nei tronchi del bosco mi rinchiuse un dì.}
\strofa
\rigo{RIT.}
\strofa
\rigo{Da quel giorno nei tronchi prigioniero son}
\rigo{e costretto a soffrire freddo e solleon.}
\strofa
\rigo{RIT.}
\strofa
\rigo{Nell’ardor della fiamma mi consumo qui}
\rigo{e col ceppo che arde brucio anch’io ogni dì.}
\strofa
\rigo{RIT.}
\strofa
\rigo{Dal tremendo supplizio convertito son}
\rigo{e per l’uomo divento fuoco caldo e buon.}
\strofa
\rigo{RIT.}
\strofa
\rigo{> Indice}
\end{multicols}

\cantotesta{29}{Alla fiera dell'est}{Fuoco di bivacco · Bans}{A. Branduardi, L. Zappa}{\bLC}{}
\begin{multicols}{2}
\rigo{\ch{Mim}{Alla}\ch{Sol}{} fiera dell'est,}
\rigo{p\ch{Re}{er} \ch{Sol}{due} soldi}
\rigo{u\ch{Re}{n} t\ch{Sol}{opolino}}
\rigo{m\ch{Si7}{io} p\ch{Mim}{adre} comprò.}
\strofa
\rigo{E\ch{Re}{} ve\ch{Sol}{nne} il gatto}
\rigo{c\ch{Re}{he} \ch{Sol}{si} mangiò il topo}
\rigo{che al mercato}
\rigo{mio padre comprò.}
\strofa
\rigo{E venne il cane}
\rigo{che morse il gatto...}
\strofa
\rigo{E venne il bastone}
\rigo{che picchiò il cane...}
\strofa
\rigo{E venne il fuoco}
\rigo{che bruciò il bastone...}
\strofa
\rigo{E venne l'acqua}
\rigo{che spense il fuoco...}
\strofa
\rigo{E venne il toro}
\rigo{che bevve l'acqua...}
\strofa
\rigo{E venne il macellaio}
\rigo{che uccise il toro...}
\strofa
\rigo{E l'angelo della morte...}
\strofa
\rigo{E infine il Signore....}
\strofa
\rigo{> Indice}
\end{multicols}

\cantotesta{30}{Il duca di Barnabò}{Fuoco di bivacco · Bans}{}{\bLC}{}
\begin{multicols}{2}
\rigo{Il conte di Barnabò}
\rigo{aveva una compagnia}
\rigo{che andava su e giù}
\rigo{oppure a mezza via.}
\rigo{Saliti ch'eran su}
\rigo{discesi ch'eran giù}
\rigo{chi stava mezza via}
\rigo{non stava ne su e ne giù.}
\strofa
\rigo{> Indice}
\end{multicols}

\cantotesta{31}{Jack è in cucina}{Fuoco di bivacco · Bans}{}{\bLC}{}
\begin{multicols}{2}
\rigo{Jack va in cucina con Tina (3 volte)}
\rigo{e suona il suo vecchio banjo}
\rigo{bom bom bom}
\rigo{fli flai fli flai flo (3 volte)}
\rigo{e suona il suo vecchio banjo.}
\strofa
\rigo{> Indice}
\end{multicols}

\cantotesta{32}{La battaglia di Magenta}{Fuoco di bivacco · Bans}{}{\bLC}{}
\begin{multicols}{2}
\rigo{Era una sera Battaglia di Magenta}
\rigo{o che piacere caricare i cavalieri}
\rigo{cavalieri caricate con una mano.}
\strofa
\rigo{(per finire:) Cavalieri carica!}
\strofa
\rigo{> Indice}
\end{multicols}

\cantotesta{33}{La macchina del capo}{Fuoco di bivacco · Bans}{}{\bLC}{}
\begin{multicols}{2}
\rigo{-La macchina del capo}
\rigo{ha un buco nella gomma,- (3 volte)}
\rigo{ripariamola col chewing-gum.}
\strofa
\rigo{(si ripete sostituendo ogni volta con i rumori)}
\strofa
\rigo{> Indice}
\end{multicols}

\cantotesta{34}{Siamo andati alla caccia del leon}{Fuoco di bivacco · Bans}{}{\bLC}{}
\begin{multicols}{2}
\rigo{\ch{Sol}{Sono} andato alla cacci\ch{Do}{a} del leon...\ch{Sol}{} pem pem}
\rigo{Sono andato alla caccia del leon... pem p\ch{Re}{em}}
\rigo{Sono andato \ch{Sol}{alla} caccia\ch{Sol7}{} (3v)\ch{Do}{} de\ch{Lam}{l} le\ch{Sol}{on..}\ch{Re}{} pe\ch{Sol}{m} pem}
\strofa
\rigo{RIT.}
\rigo{SOL               DO                 SOL (RE)}
\rigo{Singing ja ja yuppi yuppi ja (2v)}
\rigo{Singing ja j\ch{Sol}{a} yu\ch{Sol7}{ppi} ja j\ch{Do}{a} yupp\ch{Lam}{i} ja ja y\ch{Sol}{uppi} yuppi \ch{Re}{ja}      \ch{Sol}{}}
\strofa
\rigo{Taglia l'ancora la nave per partir... tu tu  RIT.}
\strofa
\rigo{Sulla riva son gli amici a salutar... bye bye  RIT.}
\strofa
\rigo{Siamo ormai nella foresta equatorial... brr brr  RIT.}
\strofa
\rigo{Un leone sta dormendo non lontan ... ron ron  RIT.}
\strofa
\rigo{Se ci vede ci divora in un boccon... gnam gnam  RIT.}
\strofa
\rigo{Cari amici non facciamolo svegliar... pss pss  RIT.}
\strofa
\rigo{> Indice}
\end{multicols}

\cantotesta{35}{Sul lago Tanganika}{Fuoco di bivacco · Bans}{}{\bEG}{}
\begin{multicols}{2}
\rigo{Sul lago Tanganika zum baibai zum baibai}
\rigo{Ci stava una piroga zum baibai zum baibai}
\rigo{Ci stava una piroga zum baibai zum baibai}
\rigo{E con il capo degli Zulù}
\strofa
\rigo{Ela mazù, ela mazù, evviva il capo degli Zulù}
\strofa
\rigo{E in quella piroga}
\rigo{Ci stavano i guerrieri (2v)}
\rigo{E con il capo degli Zulù}
\strofa
\rigo{Ela mazù, ela mazù, evviva il capo degli Zulù}
\strofa
\rigo{E i forti guerrieri}
\rigo{Andavano alla guerra (2v)}
\rigo{E con il capo degli Zulù}
\strofa
\rigo{Ela mazù, ela mazù, evviva il capo degli Zulù}
\strofa
\rigo{Andavano alla guerra}
\rigo{Perché avevano rapito (2v)}
\rigo{E la regina degli Zulù}
\strofa
\rigo{Ela mazù, ela mazù, evviva il capo degli Zulù.}
\strofa
\rigo{> Indice}
\end{multicols}

\cantotesta{36}{Un elefante}{Fuoco di bivacco · Bans}{}{\bEG}{}
\begin{multicols}{2}
\rigo{Un elefante si dondolava}
\rigo{sopra il filo di una ragnatela}
\rigo{e ritenendo la cosa interessante}
\rigo{andò a chiamare un altro elefante.}
\strofa
\rigo{Due elefanti si dondolavano...}
\rigo{Tre elefanti...}
\strofa
\rigo{> Indice}
\end{multicols}

\cantotesta{37}{Sul colle}{Fuoco di bivacco · Crescendo}{}{\bEG}{}
\begin{multicols}{2}
\rigo{Rit. Sul colle scorre il ruscello ancor,}
\rigo{nel bosco canta il cucù,}
\rigo{è sorto il sole, esplorator,}
\rigo{non indugiare più, non indugiar.}
\strofa
\rigo{È sorto il sole esplorator,}
\rigo{non indugiare più!}
\rigo{Rit.rna al tuo lavoro,}
\rigo{che grande gioia ti dà}
\rigo{diffondi attorno a te così}
\rigo{letizia e gran bontà.}
\strofa
\rigo{La buona azione non tralasciar}
\rigo{sii sempre pronto a servir;}
\rigo{ricorda sempre: nel tuo lavor}
\rigo{lealtà devi seguir.}
\strofa
\rigo{> Indice}
\end{multicols}

\cantotesta{38}{Vecchio scarpone}{Fuoco di bivacco · Crescendo}{}{\bRS}{}
\begin{multicols}{2}
\rigo{Quando la neve scende giù lieve non esitar si deve}
\rigo{prendi il tuo sacco e lesto va, dove più fioccherà.}
\rigo{Vecchio scarpone d' ogni passione vinci la tentazione}
\rigo{non abbracciar la tua Mimì, ma i tuoi fedeli sci...}
\strofa
\rigo{Si va sulla montagna}
\rigo{dove la neve il volto ci abbronzerà.}
\rigo{L'ardor che ci accompagna,}
\rigo{come una fiamma il cuor ci riscalderá.}
\rigo{Salir, sermpre salir,}
\rigo{mentre ogni valle ripeterà così... SCI SCI}
\rigo{Sciator, ripete il vento}
\rigo{con ardimento íl tuo motto sarà SCIAR}
\strofa
\rigo{Signorinella pallida e snella, lascia la tua gonnella e la pelliccia di vison,}
\rigo{mettiti i pantalon. Bon Bo Bon Bon...}
\rigo{Vecchio Scarpone d'ogni passione, vinci la tentazione,}
\rigo{non abbracciar la Mimi, ma i tuoi fedeli sci...}
\strofa
\rigo{> Indice}
\end{multicols}

\cantotesta{39}{Canzone dell'amicizia}{Fuoco di bivacco · Picco}{}{\bEG}{}
\begin{multicols}{2}
\rigo{\ch{Do}{In} \ch{Sol}{un} mondo di maschere}
\rigo{d\ch{Lam}{ove} \ch{Mim}{sembra} impossibile}
\rigo{r\ch{Fa}{ius}\ch{Do}{cire} a sconfiggere}
\rigo{t\ch{Re}{utt}\ch{Sol}{o} ciò che annienta l'uomo.}
\strofa
\rigo{Il potere, la falsità,}
\rigo{la ricchezza, l'avidità,}
\rigo{sono mostri da abbattere,}
\rigo{noi però non siamo soli.}
\strofa
\rigo{C\ch{Do}{ant}\ch{Sol}{a} con noi, batti le mani,}
\rigo{a\ch{Lam}{lzal}\ch{Mim}{e} in\ch{Fa}{} al\ch{Sol}{to,} muovile al ritmo del canto.}
\rigo{S\ch{Do}{tri}\ch{Sol}{ngi} la mano del tuo vicino}
\rigo{e\ch{Lam}{} sco\ch{Mim}{prir}\ch{Fa}{ai} \ch{Sol}{che} \ch{Do}{è} meno duro il cammino così.}
\strofa
\rigo{Ci hai promesso il tuo spirito,}
\rigo{lo sentiamo in mezzo a noi}
\rigo{e perciò possiam credere}
\rigo{che ogni cosa può cambiare.}
\strofa
\rigo{Non possiamo più assistere}
\rigo{impotenti ed attoniti}
\rigo{perché siam responsabili}
\rigo{della vita intorno a noi.}
\strofa
\rigo{RIT.}
\strofa
\rigo{> Indice}
\end{multicols}

\cantotesta{40}{Caramba}{Fuoco di bivacco · Picco}{}{\bLC}{}
\begin{multicols}{2}
\rigo{Lassù in cima al monte Nero,}
\rigo{c'è una piccola caverna}
\rigo{ci son dodici briganti,}
\rigo{al chiaror d'una lanterna.}
\strofa
\rigo{Rit. Caramba beviamo del whisky}
\rigo{caramba beviamo del gin.}
\strofa
\rigo{E tu non dar retta al cuore}
\rigo{che tutto passerà.}
\strofa
\rigo{Mentre tutti son festanti}
\rigo{uno solo resta muto:}
\rigo{ancor pieno è il suo bicchiere:}
\rigo{come mai non ha bevuto?}
\rigo{Ma non può dimenticare}
\rigo{il brigante la sua bella}
\rigo{i suoi occhi come il mare,}
\rigo{la sua bocca tanto bella.}
\strofa
\rigo{> Indice}
\end{multicols}

\cantotesta{41}{John Brown}{Fuoco di bivacco · Picco}{tradizionale (USA)}{\bEG}{}
\begin{multicols}{2}
\rigo{John Brown giace nella tomba là nel pian}
\rigo{dopo una lunga lotta contro l'oppressor;}
\rigo{John Brown giace nella tomba là nel pian,}
\rigo{ma l'anima vive ancor.}
\strofa
\rigo{Rit. Glory, glory hallelujah}
\rigo{glory, glory hallelujah,}
\rigo{glory, glory hallelujah}
\rigo{ma l'anima vive ancor.}
\strofa
\rigo{Con diciannove suoi compagni di valor,}
\rigo{dall'Est all'Ovest la Virginia conquistò;}
\rigo{con diciannove suoi compagni di valor,}
\rigo{ma l'anima vive ancor.}
\strofa
\rigo{Poi l'hanno ucciso come fosse un traditor}
\rigo{ma il traditor fu quell'uomo che l'impiccò;}
\rigo{poi l'hanno ucciso come fosse un traditor,}
\rigo{ma l'anma vive ancor.}
\strofa
\rigo{John Brown è morto ma lo schiavo è in libertà,}
\rigo{tutti fratelli, bianchi e neri siamo già.}
\strofa
\rigo{John Brown è morto, ma lo schiavo è in libertà,}
\rigo{ma l'anima vive ancor.}
\strofa
\rigo{Oh, non temere colui che il corpo ucciderà}
\rigo{se la tua anima rapir non ti potrà;}
\rigo{Oh, non temere colui che il corpo ucciderà}
\rigo{se l'anima vive ancor.}
\strofa
\rigo{Stelle dei cieli non piangete su John Brown,}
\rigo{stelle dei cieli sorridete con John Brown;}
\rigo{stelle dei cieli non piangete su John Brown,}
\rigo{la sua anima vive ancor.}
\strofa
\rigo{> Indice}
\end{multicols}

\cantotesta{42}{L'uomo di Cromagnon}{Fuoco di bivacco · Picco}{}{\bLC}{}
\begin{multicols}{2}
\rigo{Era l'alba della preistoria}
\rigo{mille e mille anni fa}
\rigo{appariva senza gloria}
\rigo{un essere simile all'orangutan;}
\rigo{aveva il corpo coperto di peli}
\rigo{l'ascia di pietra ed il baston...}
\strofa
\rigo{e volea conquistare la terra}
\rigo{era l'uomo di Cromagnon.}
\strofa
\rigo{Rit. L'uomo di CRO}
\rigo{l'uomo di MA}
\rigo{l'uomo di GNON...}
\strofa
\rigo{l'uomo di Cromagnon;}
\rigo{l'uomo di Cromagnon}
\rigo{l'era proprio un bestion}
\rigo{l'era proprio un bestion.}
\strofa
\rigo{L'uomo di Cromagnon}
\rigo{l'era proprio un bestion}
\rigo{l'era proprio un bestion.}
\strofa
\rigo{Pon pon pon para pam pon pon.}
\strofa
\rigo{Scavalcando i secoli e i mesi}
\rigo{chè la storia la sappiam:}
\rigo{lasciamo perdere gli Esquimesi}
\rigo{gli antichi Assiri e gli Egizian:}
\rigo{non parliam delle armi da guerra:}
\rigo{la spingarda ed il cannon:}
\rigo{chi vuol mai conquistare la terra}
\rigo{è ancor l'uomo di Cromagnon.}
\strofa
\rigo{> Indice}
\end{multicols}

\cantotesta{43}{La bella la va al fosso}{Fuoco di bivacco · Picco}{tradizionale}{\bRS}{}
\begin{multicols}{2}
\rigo{La bella la va al fosso.}
\rigo{Ravanel, remolass, barbabietole e spinass}
\rigo{Tre palanche al mass}
\rigo{La bella la va la fosso, al fosso a resentar}
\rigo{Sì, al fosso a resentar}
\rigo{Intant che la resenta e gh'è burlà giò l'anel}
\rigo{La alza gli occhi al cielo, la vide il ciel seren}
\rigo{La sbassa gli occhi all'onda, la vide un pescator}
\rigo{Ohi pescator dell'onda, pescatemi l'anel}
\rigo{E qul'avrai pescato un regalo ti farò}
\rigo{Andrem lassù sui monti, sui monti a far l'amor.}
\strofa
\rigo{> Indice}
\end{multicols}

\cantotesta{44}{Lazy boy}{Fuoco di bivacco · Picco}{}{\bEG}{}
\begin{multicols}{2}
\rigo{Lazy Boy}
\strofa
\rigo{C\ch{Lam}{'è} una stella bianca,}
\rigo{c\ch{Do}{he} splende su nel ciel,}
\rigo{g\ch{Lam}{uardando} solitaria}
\rigo{i\ch{Do}{} pascoli del Far West,}
\rigo{è\ch{Rem}{} la \ch{Mim}{gran} stella}
\rigo{d\ch{Fa}{el} \ch{Mim}{vecchio} Texas,}
\rigo{l\ch{Rem}{a} st\ch{Mim}{ella}\ch{Lam}{} del cow boy.}
\strofa
\rigo{Nell'Ottantasette,}
\rigo{brillando su Fort Joy,}
\rigo{vide venire al mondo}
\rigo{il piccolo Lazy Boy.}
\rigo{E qui comincia la lunga storia}
\rigo{del pallido cow-boy.}
\strofa
\rigo{RIT: DO LA-}
\rigo{Un dì la mamma gli disse va,}
\rigo{m\ch{Do}{a} r\ch{Sol}{esta}\ch{Lam}{} un bravo cow boy,}
\rigo{u\ch{Do}{n} d\ch{Lam}{ì} la tromba ti chiamerà,}
\rigo{D\ch{Do}{ani}\ch{Sol}{ele} \ch{Lam}{la} suonerà.}
\strofa
\rigo{Lazy vuol dire pigro,}
\rigo{ma il nostro Lazy Boy,}
\rigo{non era affatto pigro}
\rigo{lontano da fort Joy.}
\rigo{Ma un brutto giorno}
\rigo{conobbe Jessie,}
\rigo{conobbe la sua colt.}
\strofa
\rigo{Cominciò a sparare}
\rigo{ai sassi della via,}
\rigo{poi preferì i cavalli}
\rigo{e infine la ferrovia,}
\rigo{ai passeggieri,}
\rigo{bucò sombreri,}
\rigo{rubò pepite d'or.}
\strofa
\rigo{Tutti gli sceriffi, cercavan Lazy Boy,}
\rigo{la taglia era grossa,}
\rigo{piaceva al sergente Roy,}
\rigo{ma ad una spanna, trovò la canna}
\rigo{del pallido cow boy.}
\strofa
\rigo{Stava per sparare}
\rigo{il nostro Lazy Boy,}
\rigo{ma udì da una missione,}
\rigo{un coro di Little Boy:}
\rigo{"Oh lode, oh glory, oh alleluia}
\rigo{e il colpo non partì.}
\strofa
\rigo{C'è una stella bianca,}
\rigo{che splende su nel ciel,}
\rigo{guardando solitaria}
\rigo{i pascoli del Far West}
\rigo{e qui finisce}
\rigo{la lunga storia}
\rigo{del pallido cow boy.}
\strofa
\rigo{> Indice}
\end{multicols}

\cantotesta{45}{Me compare Giacometo}{Fuoco di bivacco · Picco}{tradizionale}{\bEG}{}
\begin{multicols}{2}
\rigo{Me compare Giacometo}
\rigo{El gh'aveva un bel galeto}
\rigo{Quando canta al verdail bevo}
\rigo{Ch'el fa propi innamorar}
\strofa
\rigo{Un bel giorno la padrona}
\rigo{Per far festa agli invitati}
\rigo{La g'ha tira il collo al gallo}
\rigo{E lo mette a cucinar}
\strofa
\rigo{Le galline tutte matte}
\rigo{Per la perdita del gallo}
\rigo{Le g'ha rotto el caponaro}
\rigo{De la rabbia che le g'ha.}
\strofa
\rigo{> Indice}
\end{multicols}

\cantotesta{46}{Non si va in cielo}{Fuoco di bivacco · Picco}{}{\bLC}{}
\begin{multicols}{2}
\rigo{\ch{Sol}{Non} si va in cielo in bicicl\ch{Re}{etta}}
\rigo{perché in cielo si va un po’ più in frett\ch{Sol}{a...}}
\rigo{Non si va in ciel\ch{Do}{o} in bicicletta}
\rigo{perché in cie\ch{Sol}{lo} si va un po’ più in frett\ch{Re}{a..}\ch{Sol}{.}   \ch{Do}{}   \ch{Sol}{}}
\strofa
\rigo{Non si va in cielo in Pininfarina}
\rigo{perché in cielo non c’è la benzina.}
\strofa
\rigo{Non si va in cielo con la vestaglia}
\rigo{perché in cielo si lavora a maglia.}
\strofa
\rigo{Non si va in cielo con il portafoglio}
\rigo{perché in cielo non c’è imbroglio.}
\strofa
\rigo{Non si va in cielo con i cardinali}
\rigo{perché in cielo sono tutti uguali.}
\strofa
\rigo{Non si va in cielo con i capitalisti}
\rigo{perché in cielo sono comunisti.}
\strofa
\rigo{Se tu vai in cielo prima di me}
\rigo{fai un buchino e tira su anche me.}
\strofa
\rigo{Se tu vai all’inferno prima di me}
\rigo{tappa tutti i buchi che io non passi più.}
\strofa
\rigo{> Indice}
\end{multicols}

\cantotesta{47}{Oleanna}{Fuoco di bivacco · Picco}{tradizionale (Norvegia)}{\bLC}{}
\begin{multicols}{2}
\rigo{\ch{Re}{Ole-Oleanna} vecchia\ch{Sol}{} terra del Far-W\ch{Re}{est}}
\rigo{f\ch{Sol}{ra} le terre\ch{Re}{} la più strana tutt\ch{La}{a} ti vorrei per me.\ch{Re}{}}
\strofa
\rigo{RIT.}
\rigo{O\ch{Re}{le-Oleanna,} Ole-Ole\ch{Sol}{anna-na,}        \ch{Re}{}}
\rigo{O\ch{Sol}{le-Ole,} Ole\ch{Re}{-Ole,} Ole-O\ch{La}{leanna-na.}\ch{Re}{}}
\strofa
\rigo{Sheeee...nengo il caballero}
\rigo{quando ad Oleanna fu}
\rigo{mise in piedi un grattacielo}
\rigo{con le fondamenta in su.}
\strofa
\rigo{RIT.}
\strofa
\rigo{Le mucche campagnole}
\rigo{tutte migrano in città}
\rigo{vi frequentano le scuole}
\rigo{dove imparano a ballar.}
\strofa
\rigo{RIT.}
\strofa
\rigo{Se vino tu vuoi bere}
\rigo{che ti faccia inebriar}
\rigo{dalle pompe di un pompiere}
\rigo{non ce l’hai che da succhiar.}
\strofa
\rigo{RIT.}
\strofa
\rigo{> Indice}
\end{multicols}

\cantotesta{48}{Sul cappello un bel fior}{Fuoco di bivacco · Picco}{}{\bEG}{}
\begin{multicols}{2}
\rigo{Voi che passare ci vedete}
\rigo{sotto la pioggia o del sol al raggio}
\rigo{qualche volta penserete}
\rigo{che ci vuole un bel coraggio}
\rigo{per andare sempre così}
\rigo{pur se ripida è la via:}
\rigo{voi V.P. non ci capite}
\strofa
\rigo{perché certo non avete...}
\strofa
\rigo{Rit. Sul cappello un bel fior}
\rigo{sulle labbra una canzon}
\rigo{un cuor gioioso e sincero}
\rigo{questo solo ci vuol}
\rigo{per un bravo Esplorator}
\rigo{per girare il mondo intero.}
\strofa
\rigo{Oh potessimo nei boschi}
\rigo{ogni nostro dì passare}
\rigo{lieti poi sotto le tende}
\rigo{alla sera riposare.}
\strofa
\rigo{Dalle vette delle Alpi}
\rigo{alle sponde del gran mare}
\rigo{con il nostro lieto canto}
\rigo{noi diremo a tutto il mondo...}
\strofa
\rigo{> Indice}
\end{multicols}

\cantotesta{49}{Ah io vorrei tornare}{Fuoco di bivacco · Calando}{}{\bRS}{}
\begin{multicols}{2}
\rigo{A\ch{Sol}{h,} i\ch{Do}{o} v\ch{Lam}{orrei} tornare anche solo per un dì}
\rigo{L\ch{Re}{ass}\ch{Sol}{ù} nella valle alpina.}
\rigo{Lassù tra gli alti abeti ed i rododendri in fior}
\rigo{Distendermi a terra e sognar.}
\strofa
\rigo{P\ch{Sol}{orta}\ch{Do}{mi} \ch{Sol}{tu} lassù, o Signor}
\rigo{D\ch{Re}{ove} meglio ti veda…}
\strofa
\rigo{P\ch{Sol}{orta}\ch{Do}{mi} \ch{Lam}{tra} il verde dei tuoi pascoli, lassù}
\rigo{P\ch{Re}{er} \ch{Sol}{non} farmi scender mai più (2 vv)}
\strofa
\rigo{Là, sotto il pino antico noi lasciammo nel partir}
\rigo{La croce del nostro altare}
\rigo{Ma sotto il pino antico con la croce là restò}
\rigo{Un poco del nostro cuor.}
\strofa
\rigo{Portami tu …}
\strofa
\rigo{E quando quest'inverno qui la neve scenderà}
\rigo{Bianca sarà la valle}
\rigo{Ma sotto il pino antico un bel giglio fiorirà}
\rigo{Il giglio dell'esplorator.}
\strofa
\rigo{> Indice}
\end{multicols}

\cantotesta{50}{Canto del ricordo}{Fuoco di bivacco · Calando}{}{\bEG}{}
\begin{multicols}{2}
\rigo{Quando tramonta il sol B.P. tu torni tra noi,}
\rigo{quando riflette la fiamma le ombre di noi}
\rigo{che cantiamo in cor; tu torni B.P.,}
\strofa
\rigo{tu torni tra noi ci guardi non visto}
\rigo{da un angolo d'ombra cantando con noi.}
\strofa
\rigo{E tu rivedi a sera i fuochi dei militar,}
\rigo{rivedi i fuochi, le tende, i cavalli, le veglie,}
\rigo{i deserti, il mar. Tu torni B.P.,}
\strofa
\rigo{tu torni a cantar; confuso nel vento}
\rigo{si perde il tuo canto col nostro che va.}
\strofa
\rigo{(... fischio... )}
\rigo{Quando tramonta il sol B.P., tu torni fra noi,}
\strofa
\rigo{quando riflette la fiamma le ombre di noi}
\rigo{che cantiamo in cor; tu torni B.P.}
\strofa
\rigo{tu torni fra noi, ritorni ogni sera}
\rigo{cantando sull'eco, dei canti d'allor.}
\strofa
\rigo{> Indice}
\end{multicols}

\cantotesta{51}{Carnet di marcia}{Fuoco di bivacco · Calando}{}{\bEG}{}
\begin{multicols}{2}
\rigo{Guardo nei campi brulli le stoppie aride}
\rigo{e nel canneto osservo levarsi un vol.}
\rigo{Mi chiedo che fanno queste cose intorno}
\rigo{è un sogno o un inganno}
\rigo{questa vita accanto a me.}
\rigo{Sei tu Signor che ti nascondi cercano Te.}
\rigo{La mia tendina chiara spicca tra gli alberi}
\rigo{nella radura erbosa declina il dì.}
\rigo{trattiene il respiro ogni cosa intorno}
\rigo{il fuoco che miro raccoglie tutto in sè.}
\rigo{Marcio con zaino in spalla per valli insolite}
\rigo{divido il pane e l'acqua con il fratel.}
\rigo{La gente che vedo mi ridà il saluto}
\rigo{le cose in cui credo son concrete accanto a me.}
\rigo{Dolci ricordi tornano}
\rigo{Ah io vorrei tornare anche solo per un dì}
\rigo{lassù nella valle alpina,}
\rigo{lassù tra gli alti abeti ed i rododendri in fior}
\rigo{distendermi a terra e sognar.}
\strofa
\rigo{Portami tu lassù, Signor}
\rigo{dove meglio ti veda!}
\rigo{Oh portami tra il verde}
\rigo{dei tuoi pascoli lassù}
\rigo{per non farmi scendere mai più.}
\rigo{Là, sotto il pino antico noi lasciammo nel partir}
\rigo{la croce del nostro altare}
\rigo{là, sotto il pino antico con la croce là restò}
\rigo{un poco del nostro cuor.}
\rigo{E quando quest'inverno qui la neve scenderà}
\rigo{bianca sarà la valle,}
\rigo{ma sopra quella croce un bel giglio fiorirà}
\rigo{il giglio dell'esplorator.}
\strofa
\rigo{> Indice}
\end{multicols}

\cantotesta{52}{Colore del sole}{Fuoco di bivacco · Calando}{}{\bRS}{}
\begin{multicols}{2}
\rigo{Colore del sol, più giallo dell'oro fino}
\rigo{la gioia che brilla negli occhi tuoi fratellino;}
\rigo{colore di un fior, il primo della mattina,}
\rigo{in te sorellina vedo rispecchiar;}
\rigo{nel gioco che fai non sai ma già stringi forte}
\rigo{segreti che sempre avrai nelle mani in ogni tua sorte;}
\rigo{nel canto che sai la voce tua chiara dice:}
\rigo{famiglia felice oggi nel mondo poi su nel ciel}
\rigo{per sempre sarà.}
\strofa
\rigo{Colore del grano che presto avrà d'or splendore,}
\rigo{la verde tua tenda fra gli alberi esploratore;}
\rigo{colore lontano di un prato che al cielo grida}
\rigo{sorella mia guida con me scoprirai...}
\strofa
\rigo{Dove il tuo sentiero ti porta in un mondo ignoto}
\rigo{che aspetta il tuo passo come di un cavaliere fidato;}
\rigo{dove il tuo pensiero dilata questa natura}
\rigo{verso l'avventura che ti farà capire te stessa}
\rigo{e amare il Signor.}
\strofa
\rigo{Colore rubino del sangue di un testimone,}
\rigo{il sangue che sgorga dal cuor tuo fratello rover;}
\rigo{color vespertino di nubi sul sol raccolte,}
\rigo{l'amore che Scolta tu impari a donar:}
\rigo{conosci quel bene che prima tu hai ricevuto}
\rigo{sai che non potrai tenerlo per te neppure un minuto;}
\rigo{conosci la gioia di spenderti in sacrificio}
\rigo{prometti servizio ad ogni fratello ovunque}
\rigo{la strada ti porterà.}
\strofa
\rigo{> Indice}
\end{multicols}

\cantotesta{53}{Faria faria oh}{Fuoco di bivacco · Calando}{Paul Doncoeur}{\bEG}{}
\begin{multicols}{2}
\rigo{Per il mondo ha camminato, faria, faria oh!}
\rigo{Monti e mari ha già varcato faria, faria oh!}
\rigo{Ma gli è dolce continuar}
\rigo{rischi e incognite a sfidar}
\rigo{faria, faria, faria, faria, faria, faria oh!}
\rigo{Non gli pesa il borsellino, faria, faria oh!}
\rigo{Ma è signor del suo destino faria, faria oh!}
\rigo{Sempre lieto e pien d'ardor canta e va l'esplorator.}
\strofa
\rigo{Faria, faria, faria, faria, faria, faria oh!}
\rigo{Col fagotto e col bastone, faria, faria oh!}
\rigo{Parte senza colazione, faria, faria oh!}
\rigo{Lunga strada far dovrà molti amici troverà.}
\strofa
\rigo{Faria, faria, faria, faria, faria, faria oh!}
\rigo{Parte e sogna cose belle, faria, faria oh!}
\rigo{Albe, fiori, prati e stelle, faria, faria oh!}
\rigo{Quando Iddio lo chiamerà sorridendo tornerà}
\rigo{Faria, faria, faria, faria, faria, faria oh!}
\strofa
\rigo{> Indice}
\end{multicols}

\cantotesta{54}{Guarda lontano}{Fuoco di bivacco · Calando}{}{\bEG}{}
\begin{multicols}{2}
\rigo{Sono gesti ispirati}
\rigo{Trasparenti colori:}
\rigo{C' è un pittore invisibile}
\rigo{che dipinge un nuovo mondo che}
\rigo{prende forma già davanti a noi;}
\rigo{la sua tela è la terra}
\rigo{la cornice il cielo blu.}
\rigo{Sono note di musica}
\rigo{nuovi suonì un pò magici:}
\rigo{un musicista invisibile}
\rigo{sta orchestrando una canzone che}
\rigo{si fa strada proprio grazie a lui;}
\rigo{toccherà nuove corde che}
\rigo{suoneranno dentro noi.}
\rigo{Uomini uomini d' oggi}
\rigo{che guardano lontano uomini}
\rigo{presi da un vento e portati in una mano.}
\rigo{E’ disegno insondabile}
\rigo{che da un senso di vertigine:}
\rigo{chi di noi è ripetibile?}
\rigo{Siamo tutti strumenti}
\rigo{che lascian chiaro un segno dietrodi sé}
\rigo{siamo fatti per vivere.}
\strofa
\rigo{> Indice}
\end{multicols}

\cantotesta{55}{Kamaludu}{Fuoco di bivacco · Calando}{}{\bEG}{}
\begin{multicols}{2}
\rigo{Kamaludu, kamaludu,}
\rigo{ludu kamaludu, uandaisii}
\rigo{kamaludu, kamaludu uandaisii}
\rigo{Behbel ina, behbel ina,}
\rigo{ina behbel ina, uandaisi}
\rigo{behbel ina, behbel ina uandaisii}
\rigo{Matulefa, matulefa,}
\rigo{efa matulefa, uandaisi}
\rigo{matulefa, matulefa uandaisi.}
\strofa
\rigo{> Indice}
\end{multicols}

\cantotesta{56}{La mula de Parenzo}{Fuoco di bivacco · Calando}{}{\bEG}{}
\begin{multicols}{2}
\rigo{La mula de Parenzo}
\rigo{L'ha messo su bottega}
\rigo{E tutto la vendeva}
\rigo{E tutto la vendeva}
\strofa
\rigo{La mula de Parenzo}
\rigo{L'ha messo su bottega}
\rigo{E tutto la vendeva}
\rigo{Föra che 'l baccalà}
\strofa
\rigo{Perchè non m'ami più}
\strofa
\rigo{La me morosa vegia}
\rigo{La tengo di riserva}
\rigo{Ma quando spunta l'erba}
\rigo{La mando a pascolar}
\strofa
\rigo{La mando a pascolare}
\rigo{Insieme alle caprette}
\rigo{L'amor con le servette}
\rigo{Non lo farò mai più.}
\strofa
\rigo{Tutti mi dicono bionda}
\rigo{Ma io bionda non sono}
\rigo{Porto i capelli neri}
\rigo{Neri come il carbon.}
\strofa
\rigo{Se il mare fosse cacio}
\rigo{E i monti di polenta}
\rigo{Ohi, mamma che pociate}
\rigo{Polenta e baccalà.}
\strofa
\rigo{> Indice}
\end{multicols}

\cantotesta{57}{Pim pam}{Fuoco di bivacco · Calando}{}{\bEG}{}
\begin{multicols}{2}
\rigo{Nella casa sulla montagna}
\rigo{un camino grande grande sta,}
\rigo{nel camino grande grande grande}
\rigo{un gran fuoco fuoco fuoco va.}
\strofa
\rigo{RIT: Percio' pim pam, le scarpe pim pam,}
\rigo{di notte fan sul sentiero di pietre grosse,}
\rigo{pim pam, le scarpe pim pam,}
\rigo{di notte fan sul sentiero cosi'.}
\strofa
\rigo{Nella casa sulla montagna}
\rigo{un signore grande grande sta,}
\rigo{nella casa verde viola e bianca}
\rigo{tante sedie rosse e gialle ha.}
\strofa
\rigo{RIT:...}
\strofa
\rigo{Nella casa sulla montagna}
\rigo{una sedia a tutti tutti da,}
\rigo{e ciascuno togliera' le scarpe}
\rigo{tutti insieme poi si danzera'.}
\strofa
\rigo{> Indice}
\end{multicols}

\cantotesta{58}{Terra della betulla}{Fuoco di bivacco · Calando}{}{\bEG}{}
\begin{multicols}{2}
\rigo{\ch{Rem}{Terra} di betull\ch{Fa}{a} casa \ch{Rem}{del} castoro,    \ch{Fa}{}}
\rigo{l\ch{Sib}{à} dove errando \ch{Fa}{va} il lupo \ch{Solm}{ancora.}      \ch{Rem}{}}
\rigo{V\ch{Sib}{oglio} tornare a\ch{Fa}{ncor} al mio bel\ch{Solm}{} lago blu\ch{Sib}{.}               \ch{Rem}{}}
\strofa
\rigo{RIT.}
\rigo{B\ch{Rem}{umbadiaidi,} bumbadia\ch{Lam}{idi,} bumbadiaidi bum\ch{Solm}{}                 \ch{Rem}{}}
\rigo{B\ch{Rem}{umbadiaidi,} bumbadia\ch{Lam}{idi,} bumbadiaidi bum\ch{Solm}{.}                \ch{Rem}{}}
\strofa
\rigo{La mia canoa scivola leggera}
\rigo{sulle lucenti vie del grande fiume.}
\rigo{Voglio tornare ancor al mio bel lago blu.}
\strofa
\rigo{RIT.}
\strofa
\rigo{La tra gli abeti la luna appare,}
\rigo{il tuo bel viso , mamma, vedo ancora.}
\rigo{Voglio tornare ancora al mio bel lago blu.}
\strofa
\rigo{RIT.}
\strofa
\rigo{> Indice}
\end{multicols}

\cantotesta{59}{Vento della sera}{Fuoco di bivacco · Calando}{}{\bEG}{}
\begin{multicols}{2}
\rigo{Rit. Vento della sera tiepida e leggera,}
\rigo{in quest'atmosfera di serenità;}
\rigo{tutti attorno al fuoco, riposiamo un poco,}
\rigo{e cantiamo assieme le vecchie canzon.}
\strofa
\rigo{Cantiamo anche per voi fratelli lontan,}
\rigo{un giorno voi partendo ci deste la man;}
\rigo{qui attorno vi vediamo, voi siete ancora qui,}
\rigo{insieme a voi cantiamo, come cantammo un dì.}
\strofa
\rigo{Un giorno voi partiste lasciandoci qua,}
\rigo{il fuoco si spegneva nell'oscurità.}
\strofa
\rigo{Sentite: noi vi amiamo come vi amammo un dì:}
\rigo{tornate vi preghiamo, noi vi aspettiamo qui.}
\strofa
\rigo{> Indice}
\end{multicols}

\cantotesta{60}{Voglio girare il mondo}{Fuoco di bivacco · Calando}{I girasoli}{\bRS}{}
\begin{multicols}{2}
\rigo{Voglio girare il mondo}
\rigo{anche se la gente ride,}
\rigo{anche se tutti quanti}
\rigo{mi consigliano di restare.}
\strofa
\rigo{Rit: E voglio andarmene via}
\rigo{proprio perchè secondo loro}
\rigo{dormire sotto le stelle,}
\rigo{vuol dire solo sentir freddo.}
\strofa
\rigo{Voglio girare il mondo anche se sul mio cammino}
\rigo{incontrerò la pioggia che scendendo mi bagnerà}
\rigo{anche con l’acqua negli occhi}
\rigo{potrò veder tutte le cose}
\rigo{che in una stanza all’asciutto}
\rigo{io non potrei sognare mai.}
\strofa
\rigo{Rit: E voglio andarmene via}
\rigo{proprio perchè secondo loro}
\rigo{dormire sotto le stelle,}
\rigo{vuol dire solo sentir freddo.}
\strofa
\rigo{> Indice}
\end{multicols}

\cantotesta{61}{Il falco}{Fuoco di bivacco · Chiusura}{}{\bEG}{}
\begin{multicols}{2}
\rigo{Un fal\ch{Re}{co} volava nel cielo\ch{Sol}{} un mattino}
\rigo{r\ch{La}{icordo} quel tempo quando ero bambin\ch{Re}{o}}
\rigo{io lo seguivo nel ros\ch{Sol}{so} tramonto}
\rigo{dall’\ch{La}{alto} di un monte vedevo il \ch{Re}{suo} mondo.}
\strofa
\rigo{RIT.}
\rigo{E allora \ch{Sol}{eha,} eha eha ...}
\rigo{e allora \ch{Re}{eha,} eha}
\rigo{eha, eha, eha, eha, eha, \ch{Sol}{eha} eh.     \ch{Re}{}}
\strofa
\rigo{Fiumi mari e boschi senza confine}
\rigo{i chiari orizzonti e le verdi colline}
\rigo{e quando partivo per un lungo sentiero}
\rigo{partivo ragazzo e tornavo guerriero.}
\strofa
\rigo{RIT.}
\strofa
\rigo{Le tende rosse vicino al torrente}
\rigo{la vita felice tra la mia gente}
\rigo{e quando il mio arco colpiva lontano}
\rigo{sentivo l’orgoglio di essere indiano.}
\strofa
\rigo{RIT.}
\strofa
\rigo{Fiumi mari e boschi mossi dal vento}
\rigo{luna su luna i miei capelli d’argento}
\rigo{e quando era l’ora dell’ultimo sonno}
\rigo{partivo dal campo per non farvi ritorno.}
\strofa
\rigo{RIT.}
\strofa
\rigo{Un falco volava nel cielo un mattino}
\rigo{e verso il sole mi indicava il cammino}
\rigo{un falco che un giorno era stato colpito}
\rigo{ma no, non è morto, era solo ferito.}
\strofa
\rigo{RIT.}
\strofa
\rigo{> Indice}
\end{multicols}

\cantotesta{62}{Signor, fra le tende schierati}{Fuoco di bivacco · Ultimi}{}{\bEG}{}
\begin{multicols}{2}
\rigo{\ch{Sol}{Sign}\ch{Do}{or} tra le tende sch\ch{Sol}{ierati}}
\rigo{per s\ch{Do}{alutar}\ch{Sol}{} il dì che muor \ch{Re}{}       \ch{Re7}{}}
\rigo{l\ch{Sol}{e} no\ch{Do}{te} di canti accorat\ch{Sol}{i}}
\rigo{levia\ch{Do}{mo} \ch{Sol}{a} te, calde d'am\ch{Re}{or.}    \ch{Re7}{}}
\strofa
\rigo{\ch{Sol}{Ascolta} tu \ch{Do}{l'umil} preghier\ch{Sol}{a}}
\rigo{che d'aspro s\ch{Do}{uol} s'ode innalzar \ch{Re}{}}
\rigo{a\ch{Sol}{} te cui\ch{Mim}{} mancava, la sera,  \ch{Re}{}}
\rigo{un te\ch{Do}{tto} ancor per ripos\ch{Re}{ar.}        \ch{Sol}{}}
\rigo{C\ch{Sol}{hiedon} sol tutti i nostri cu\ch{Re}{ori}}
\rigo{a\ch{Re7}{} te sempre meglio servir\ch{Sol}{,}}
\rigo{g\ch{Do}{enufle}\ch{Sol}{tton} qu\ch{Re}{i} nel pian \ch{Sol}{}}
\rigo{i tuoi esplor\ch{Do}{ator.}      \ch{Dom}{}}
\rigo{Tu dal ci\ch{Sol}{el} bene\ch{Re}{dici} o Signor. \ch{Do}{}       \ch{Dom}{}       \ch{Sol}{}}
\strofa
\rigo{> Indice}
\end{multicols}

\cantotesta{63}{Era una notte che pioveva}{Canti alpini}{tradizionale}{\bRS}{}
\begin{multicols}{2}
\rigo{Era una notte che pioveva}
\rigo{E che tirava un forte vento}
\rigo{Immaginatevi che grande tormento}
\rigo{Per un alpino che stava a vegliar (bis)}
\strofa
\rigo{A mezzanotte arriva il cambio}
\rigo{Accompagnato dal capo posto}
\rigo{Ohi sentinella ritorna al tuo posto}
\rigo{Sotto la tenda a ri a riposar (bis)}
\strofa
\rigo{Quando fui stato sotto la tenda}
\rigo{Sentii un rumore giù per la valle}
\rigo{Sentivo l'acqua giù per le spalle}
\rigo{Sentivo i sassi a ro a rotolar (bis)}
\strofa
\rigo{Quando fui stato addormentato}
\rigo{Sognavo di esser con la mia bella}
\rigo{Invece ero di sentinella}
\rigo{Fare la guardia allo, allo stranier (bis)}
\strofa
\rigo{> Indice}
\end{multicols}

\cantotesta{64}{Quel mazzolin di fiori}{Canti alpini}{tradizionale}{\bRS}{}
\begin{multicols}{2}
\rigo{Quel mazzolin di fiori che vien dalla montagna …}
\rigo{E bada ben che non si bagna}
\rigo{Che lo voglio regalar}
\strofa
\rigo{Lo voglio regalare perché l'è un bel mazzetto …}
\rigo{Lo voglio dare al mio moretto}
\rigo{Questa sera quando vien}
\strofa
\rigo{Stasera quando viene sarà una brutta sera …}
\rigo{E lui sabato sera,}
\rigo{Lui non è venuto a me}
\strofa
\rigo{Non è venuto a me, l'è andai dalla Rosina …}
\rigo{E perché mi son poverina}
\rigo{Mi fa pianger, sospirar.}
\strofa
\rigo{Fa pianger, sospirare sul letto dei lamenti …}
\rigo{E cosa mai dirà la gente}
\rigo{Cosa mai dirà di me}
\strofa
\rigo{Dirà che son tradita, tradita nell' amore …}
\rigo{E perché mì, me piange il cuore}
\rigo{E per sempre piangerà}
\strofa
\rigo{> Indice}
\end{multicols}

\cantotesta{65}{Signore delle cime}{Canti alpini}{B. De Marzi}{\bRS}{}
\begin{multicols}{2}
\rigo{Dio del cielo,}
\rigo{signore delle cime}
\rigo{un nostro amico}
\rigo{,hai chiesto alla montagna.}
\rigo{Ma ti preghiamo,}
\rigo{ma ti preghiamo}
\rigo{su nel paradiso,}
\rigo{su nel paradiso}
\rigo{lascialo andare}
\rigo{,per le tue montagne.}
\rigo{Santa Maria,}
\rigo{signora della neve}
\rigo{copri col bianco,}
\rigo{soffice mantello}
\rigo{il nostro amico}
\rigo{il nostro fratello.}
\rigo{Su nel paradiso,}
\rigo{su nel paradiso}
\rigo{lascialo andare,}
\rigo{per le tue montagne.}
\strofa
\rigo{> Indice}
\end{multicols}

\cantotesta{66}{Sul Monte Nero}{Canti alpini}{}{\bRS}{}
\begin{multicols}{2}
\rigo{Lassù sul monte nero}
\rigo{c'è una piccola taverna;}
\rigo{ci son dodici briganti}
\rigo{al chiaror di una lanterna.}
\strofa
\rigo{Rit.: Caramba beviamo del whisky}
\rigo{caramba beviamo del gin !}
\rigo{E tu non dar retta al cuore}
\rigo{che tutto passa e và...}
\strofa
\rigo{Mentre gli altri son festanti}
\rigo{uno solo resta muto}
\rigo{ha il bicchiere ancora pieno}
\rigo{come mai non ha bevuto?}
\strofa
\rigo{Rit.:}
\strofa
\rigo{Ma non può dimenticare}
\rigo{il brigante la sua bella}
\rigo{i sui occhi luminosi}
\rigo{la sua bocca tanto bella.}
\strofa
\rigo{Rit.:}
\strofa
\rigo{Mentre stava spigolando}
\rigo{in un campo di frumento}
\rigo{è passato un aeroplano}
\rigo{che compì un mitragliamento.}
\strofa
\rigo{Rit.:}
\strofa
\rigo{Or la bella par che dorma}
\rigo{con la spiga ancora in mano}
\rigo{par che pensi al suo brigante}
\rigo{che la piange da lontano.}
\strofa
\rigo{Rit.:}
\strofa
\rigo{> Indice}
\end{multicols}

\cantotesta{67}{Ta-pum}{Canti alpini}{tradizionale}{\bRS}{}
\begin{multicols}{2}
\rigo{Ho lasciato la mamma mia,}
\rigo{l'ho lasciata per fare il soldà.}
\strofa
\rigo{Rit. Tapum, tapum, tapum !}
\rigo{Venti giorni su l'Ortigara,}
\rigo{senza cambio per dismontar.}
\strofa
\rigo{Quando portano la pagnotta,}
\rigo{il cecchino comincia a sparar.}
\strofa
\rigo{E domani si andrà all'assalto:}
\rigo{soldatino, non farti ammazzar.}
\strofa
\rigo{Quando poi si discende al piano,}
\rigo{battaglione non ha più soldà.}
\strofa
\rigo{Nella valle c'è un cimitero,}
\rigo{cimitero dei nostri soldà.}
\strofa
\rigo{Cimitero di noi soldà,}
\rigo{forse un giorno ti vengo a trovar.}
\strofa
\rigo{Battaglion di tutti i morti,}
\rigo{noi giuriamo l'Italia salvar.}
\strofa
\rigo{> Indice}
\end{multicols}

\cantotesta{68}{Acqua siamo noi}{Canti liturgici · Ingresso}{}{}{}
\begin{multicols}{2}
\rigo{> Indice}
\end{multicols}

\cantotesta{69}{Camminerò}{Canti liturgici · Ingresso}{}{}{}
\begin{multicols}{2}
\rigo{RIT.}
\rigo{C\ch{Re}{amminerò,} camminerò\ch{Sim}{}}
\rigo{s\ch{Sol}{ulla} Tua strada, Signor.\ch{La}{}}
\rigo{D\ch{Re}{ammi} la mano, voglio restar\ch{Sim}{}}
\rigo{per s\ch{Sol}{empre} insie\ch{La}{me} a Te.       \ch{Re}{}}
\strofa
\rigo{Q\ch{Re}{uando} ero solo, solo e stanco d\ch{Sim}{el} mondo,}
\rigo{q\ch{Sol}{uando} non c'era l'Amor, \ch{La}{}}
\rigo{t\ch{Re}{ante} persone vidi i\ch{Sim}{ntorno} a me,}
\rigo{s\ch{Sol}{entivo} cant\ch{La}{are} così.  \ch{Re}{}}
\strofa
\rigo{RIT.}
\strofa
\rigo{Io non capivo ma rimasi a sentire}
\rigo{quando il Signore mi parlò:}
\rigo{Lui mi chiamava, chiamava anche me}
\rigo{e la mia risposta s'alzò.}
\strofa
\rigo{RIT.}
\strofa
\rigo{Or non importa se uno ride di me,}
\rigo{lui certamente non sa}
\rigo{del gran regalo che ebbi quel dì}
\rigo{che dissi al Signore così.}
\strofa
\rigo{RIT.}
\strofa
\rigo{A volte sono triste, ma mi guardo intorno,}
\rigo{scopro il mondo e l'amor;}
\rigo{son questi i doni che Lui fa a me}
\rigo{felice ritorno a cantar.}
\strofa
\rigo{RIT.}
\strofa
\rigo{> Indice}
\end{multicols}

\cantotesta{70}{Cantiamo a te}{Canti liturgici · Ingresso}{}{}{}
\begin{multicols}{2}
\rigo{Cantiamo a Te, Signore della vita:}
\rigo{il nome tuo è grande sulla terra}
\rigo{tutto parla di Te e canta la tua gloria.}
\rigo{Grande Tu sei e compi meraviglie: Tu sei Dio.}
\strofa
\rigo{Cantiamo Te, Signore Gesù Cristo:}
\rigo{Figlio di Dio venuto sulla terra,}
\rigo{fatto uomo per noi nel grembo di Maria.}
\rigo{Dolce Gesù risorto dalla Morte sei con noi.}
\strofa
\rigo{Cantiamo Te, Amore senza fine:}
\rigo{Tu che sei Dio lo Spirito del Padre}
\rigo{Brani e testi contenuti : 1145}
\rigo{vivi dentro di noi e guidi i nostri passi.}
\rigo{Accendi in noi il fuoco dell’eterna carità.}
\strofa
\rigo{Cantiamo la nostra lode}
\rigo{A Cristo Signore,}
\rigo{vivente nella Chiesa.}
\strofa
\begin{ritornello}
\rigo{Cantiamo la nostra lode,}
\rigo{cantiamo il nostro amen!}
\rigo{Alleluia! Alleluia!}
\strofa
\end{ritornello}
\rigo{A Cristo pastore,}
\rigo{presente nella Chiesa.}
\strofa
\rigo{Cantiamo le meraviglie}
\rigo{Ho visto una folla immensa:}
\rigo{cantavano un canto ineffabile. Alleluia!}
\strofa
\begin{ritornello}
\rigo{Cantiamo le meraviglie}
\rigo{del Signore dell’universo.}
\rigo{Le genti innanzi a lui verranno}
\rigo{inneggiando per l’eternità.}
\strofa
\end{ritornello}
\rigo{Ho visto volti di gioia,}
\rigo{ho udito arpe celesti. Alleluia!}
\strofa
\rigo{L’agnello che vinto la morte,}
\rigo{nel sangue ha portato vittoria. Alleluia!}
\strofa
\rigo{Cantiamo l’inizio}
\rigo{= E non videro che Gesù}
\strofa
\begin{ritornello}
\rigo{Cantico dei redenti}
\rigo{Il Signore è la mia salvezza}
\rigo{e con lui non temo più,}
\rigo{perché ho nel cuore la certezza,}
\rigo{la salvezza è qui con me.}
\strofa
\end{ritornello}
\rigo{Ti lodo Signore perché}
\rigo{un giorno eri lontano da me:}
\rigo{ora invece sei tornato}
\rigo{e mi hai preso con te.}
\strofa
\rigo{Berrete con gioia alle fonti,}
\rigo{alle fonti della salvezza}
\rigo{e quel giorno voi direte:}
\rigo{lodate il Signore invocate il suo nome.}
\strofa
\rigo{Fate conoscere ai popoli}
\rigo{tutto quello che lui ha compiuto}
\rigo{e ricordino per sempre,}
\rigo{ricordino sempre che il suo nome è grande.}
\strofa
\rigo{Cantate a chi ha fatto grandezze}
\rigo{e si è fatto sapere nel mondo;}
\rigo{grida forte la tua gioia, abitante di Sion,}
\rigo{perché grande con te è il Signore.}
\strofa
\rigo{Brani e testi contenuti : 1145}
\rigo{Cantico della vergine}
\rigo{= Magnificat Magnifica il Signore anima mia}
\strofa
\rigo{Cantico delle creature Benediciamo il Signore}
\rigo{= Benediciamo il Signore}
\strofa
\rigo{Cantico delle Creature laudato si’}
\rigo{Laudato si’ mi’ Signore}
\rigo{per frate sole, sora luna,}
\rigo{frate vento, il cielo, le stelle}
\rigo{per sora acqua, frate focu.}
\strofa
\begin{ritornello}
\rigo{Laudato si’ mi’ Signore}
\rigo{per la terra e le tue creature}
\strofa
\end{ritornello}
\rigo{Laudato si’ mi’ Signore}
\rigo{quello che porta la tua pace:}
\rigo{e saprà perdonare,}
\rigo{per il tuo amore saprà amare.}
\strofa
\rigo{Laudato si’ mi’ Signore}
\rigo{per sora morte corporale}
\rigo{dalla quale homo vivente}
\rigo{non potrà mai, mai scappare}
\strofa
\rigo{Cantico di liberazione}
\rigo{= Voglio cantare}
\strofa
\rigo{Cantico di nozze}
\rigo{Tu che inventi per ognuno di noi}
\rigo{un cammino che porta verso te.}
\rigo{Tu che chiami ognuno per il suo nome, Tu,}
\rigo{oggi ci chiami insieme.}
\rigo{Tu che hai in mente per ognuno di noi}
\rigo{una strada bellissima con te,}
\rigo{Tu stavolta parli di un’avventura che,}
\rigo{che si percorre insieme.}
\strofa
\begin{ritornello}
\rigo{Rendi agile il vascello dell’amore nostro}
\rigo{docile al soffio del tuo spirito.}
\rigo{E per gli immaginabili mari conducilo,}
\rigo{quelli che tu sai.}
\rigo{Rendi agile il vascello dell’amore nostro}
\rigo{che tocchi lidi che tu solo sai}
\rigo{e la scia della rotta disegni un’onda:}
\rigo{quella che tu vuoi.}
\strofa
\end{ritornello}
\rigo{Dacci forza perché ognuno di noi}
\rigo{con la vita dispieghi la tua idea.}
\rigo{Dacci forza per svelare il pensiero tuo}
\rigo{su questa vita insieme.}
\strofa
\begin{ritornello}
\rigo{Cantico di Zaccaria}
\rigo{Sia benedetto il Signore, Dio d’Abramo,}
\rigo{perché ha visitato e redento il suo popolo,}
\rigo{come aveva promesso, come aveva promesso.}
\strofa
\end{ritornello}
\rigo{È la salvezza da tutti i nostri nemici}
\rigo{e dalla mani di quanti ci odiano.}
\rigo{Ha avuto misericordia di noi,}
\rigo{si è ricordato del suo patto santo.}
\rigo{Gioisci popolo fedele, perché da quel giorno}
\rigo{tu non sei più schiavo}
\rigo{un’alba nuova è sorta per te,}
\rigo{un’alba nuova è sorta per te.}
\rigo{Brani e testi contenuti : 1145}
\strofa
\rigo{Quando sei stanco, quando non sai più cosa fare,}
\rigo{quando sei triste e quando sei tanto solo}
\rigo{e perfino quando la vita ti lascia}
\rigo{ricordati che è stata detta una parola di vita.}
\strofa
\rigo{Gioisci popolo fedele, perché da quel giorno}
\rigo{tu non sei più schiavo}
\rigo{un’alba nuova è sorta per te,}
\rigo{un’alba nuova è sorta per te.}
\end{multicols}

\cantotesta{71}{Che gioia ci hai dato}{Canti liturgici · Ingresso}{}{}{}
\begin{multicols}{2}
\rigo{Che gioia ci hai dato, Signore del cielo,}
\rigo{Signore del grande universo!}
\rigo{Che gioia ci hai dato vestito di luce,}
\rigo{vestito di gloria infinita,}
\rigo{vestito di gloria infinita!}
\strofa
\rigo{Vederti risorto, vederti Signore,}
\rigo{il cuore sta per impazzire!}
\rigo{Tu sei ritornato, tu sei qui tra noi,}
\rigo{e adesso ti avremo per sempre,}
\rigo{e adesso ti avremo per sempre.}
\strofa
\rigo{Chi cercate, donne, quaggiù?}
\rigo{Chi cercate, donne, quaggiù?}
\rigo{Quello che era morto non è qui:}
\rigo{è risorto sì come aveva detto anche a voi,}
\rigo{voi gridate a tutti che è risorto lui,}
\rigo{a tutti che è risorto lui.}
\strofa
\rigo{Tu hai vinto il mondo, Gesù!}
\rigo{Tu hai vinto il mondo, Gesù!}
\rigo{Liberiamo la felicità}
\rigo{e la morte, no, non esiste più, l’hai vinta tu,}
\rigo{hai salvato tutti noi, uomini con te,}
\rigo{tutti noi, uomini con te.}
\end{multicols}

\cantotesta{72}{Emmanuel}{Canti liturgici · Ingresso}{}{}{}
\begin{multicols}{2}
\rigo{\ch{Mim}{Come} in mezzo al mare brilla un \ch{Re}{riflesso} del so\ch{Mim}{le}}
\rigo{come tra le nuvole un\ch{Re}{} buco mostra il\ch{Mim}{} cielo azzurro}
\rigo{c\ch{Sol}{ome} sopra un pi\ch{Re}{cco} c'è una stella \ch{Do}{alpina} che     \ch{Mim}{}}
\rigo{è velluto bia\ch{Re}{nco} tra le rocc\ch{Do}{e} di lassù.    \ch{Mim}{}}
\strofa
\rigo{Tu vicino alla mia anima e i miei figli in te}
\rigo{la mia strada verso la luce è l'avventura in Te,}
\rigo{nuova vita, nuovo amore in questo mondo che}
\rigo{ha dimenticato il tempo della libertà.}
\strofa
\rigo{Emmanuel, Emmanuel, Emmanuel, Dio con noi}
\rigo{nel mio amore, nel nostro amore}
\rigo{c'è l'amore di Dio con noi}
\rigo{nella nostra vita Lui ha preso il posto suo,}
\rigo{Primo posto al primo posto è venuto Lui.}
\strofa
\rigo{Una Chiesa non è soltanto una casa di mattoni}
\rigo{una Chiesa non è soltanto un altare colorato,}
\rigo{una Chiesa è la realtà che nasce tra noi}
\rigo{se ci amiamo se ci amiamo con l'amore di Dio.}
\strofa
\rigo{Cambierà cambierà cambierà l'umanità}
\rigo{cambierà cambierà cambierà l'umanità}
\rigo{nella terra è seminato e un mondo nascerà}
\rigo{nascerà nascerà una nuova umanità}
\strofa
\rigo{> Indice}
\end{multicols}

\cantotesta{73}{Grandi cose}{Canti liturgici · Ingresso}{}{}{}
\begin{multicols}{2}
\rigo{\ch{Do}{Grandi} cos\ch{Sol}{e} ha fatto il Si\ch{Rem}{gnore} per no\ch{Lam}{i}}
\rigo{h\ch{Do}{a} fatto germog\ch{Fa}{liare} i fiori tra l\ch{Do}{e} rocce.       \ch{Sol}{}}
\rigo{G\ch{Do}{randi} cose\ch{Sol}{} ha fatto il Sig\ch{Rem}{nore} per noi\ch{Lam}{}}
\rigo{c\ch{Do}{i} ha riportati\ch{Fa}{} liberi all\ch{Do}{a} nostra terra.\ch{Sol}{}}
\rigo{LA-                     RE-     7                 MI-7}
\rigo{Ed ora possiamo cantare, possiamo gridare}
\rigo{l'amore che Dio h\ch{Fa}{a} versato su noi.      \ch{Sol}{}}
\strofa
\rigo{\ch{Do}{Tu} che sai\ch{Sol}{} strapp\ch{Rem}{are} dalla morte     \ch{Lam}{}}
\rigo{h\ch{Do}{ai} solleva\ch{Fa}{to} il nostro viso d\ch{Do}{alla} polvere.  \ch{Sol}{}}
\rigo{T\ch{Do}{u} che hai \ch{Sol}{sentito}\ch{Rem}{} il nostro pianto   \ch{Lam}{}}
\rigo{n\ch{Do}{el} nostro cuor\ch{Fa}{e} hai messo un seme di \ch{Do}{felicità.}  \ch{Sol}{}  RIT.}
\strofa
\rigo{> Indice}
\end{multicols}

\cantotesta{74}{Vivere la vita}{Canti liturgici · Ingresso}{}{}{}
\begin{multicols}{2}
\rigo{Vivere la vita con le gioie e coi dolori di ogni giorno}
\rigo{è quello che Dio vuole da te.}
\strofa
\rigo{Vivere la vita e inabissarti nell’amore è il tuo destino}
\rigo{è quello che Dio vuole da te.}
\strofa
\rigo{Fare insieme agli altri la tua strada verso Lui correre con i fratelli tuoi...}
\rigo{Scoprirai allora il cielo dentro di te, una scia di luce lascerai.}
\strofa
\rigo{Vivere la vita è l’avventura più stupenda dell’amore,}
\rigo{è quello che Dio vuole da te.}
\strofa
\rigo{Vivere la vita e generare ogni momento il Paradiso}
\rigo{è quello che Dio vuole da te.}
\strofa
\rigo{Vivere perché ritorni al mondo l’unità perché dio sta nei fratelli tuoi...}
\rigo{Scoprirai allora il cielo dentro di te,}
\rigo{una scia di luce lascerai.}
\strofa
\rigo{(...musica)}
\strofa
\rigo{Vivere perché ritorni al mondo l’unità perché dio sta nei fratelli tuoi...}
\rigo{Scoprirai allora il cielo dentro di te, una scia di luce lascerai.}
\rigo{una scia di luce lascerai.}
\strofa
\rigo{> Indice}
\end{multicols}

\cantotesta{75}{Vocazione}{Canti liturgici · Ingresso}{P. Sequeri}{}{}
\begin{multicols}{2}
\rigo{\ch{Do}{Era} un giorno come \ch{Sol}{tanti} altri \ch{Fa}{}}
\rigo{e quel gi\ch{Sol}{orno} Lui pas\ch{Do}{sò.}       \ch{Fa}{}   \ch{Do}{}   \ch{Sol}{}}
\rigo{E\ch{Do}{ra} un uomo come tutt\ch{Sol}{i} gli altri \ch{Fa}{}}
\rigo{e passand\ch{Sol}{o} mi chiamò.\ch{Do}{}          \ch{Fa}{}   \ch{Do}{}   \ch{Mi}{}}
\rigo{C\ch{Lam}{ome} lo sapesse \ch{Mim}{che} il mio nome         \ch{Fa}{}}
\rigo{era p\ch{Sol}{roprio} quello,}
\rigo{c\ch{Do}{ome} mai vedess\ch{Sol}{e} proprio me        \ch{Fa}{}}
\rigo{nella Sua vit\ch{Sol}{a,} non lo so\ch{Do}{.}         \ch{Fa}{}   \ch{Do}{}   \ch{Sol}{}}
\rigo{E\ch{Do}{ra} un giorno come t\ch{Sol}{anti} altri  \ch{Fa}{}}
\rigo{e quel gi\ch{Sol}{orno} mi chiamò. \ch{Do}{}           \ch{Fa}{}   \ch{Do}{}   \ch{Mi}{}}
\strofa
\rigo{RIT.}
\rigo{T\ch{Lam}{u,} Dio,\ch{Mim}{} che co\ch{Fa}{nosci} il nome mio  \ch{Sol}{}}
\rigo{f\ch{Lam}{a'} c\ch{Mim}{he} asco\ch{Fa}{ltando} la Tua v\ch{Sol}{oce}}
\rigo{i\ch{Do}{o} ricordi \ch{Sol}{dove} porta l\ch{Fa}{a} mia strada       \ch{Sol}{}}
\rigo{n\ch{Do}{ella} vita \ch{Sol}{all'incontro}\ch{Fa}{} con Te.       \ch{Do}{}}
\strofa
\rigo{Era un'alba triste e senza vita}
\rigo{e qualcuno mi chiamò.}
\rigo{Era un uomo come tanti altri}
\rigo{ma la voce, quella no;}
\rigo{quante volte un uomo}
\rigo{con il nome giusto mi ha chiamato,}
\rigo{una volta sola l'ho sentito}
\rigo{pronunciare con amore.}
\rigo{Era un uomo come nessun altro}
\rigo{e quel giorno mi chiamò.}
\strofa
\rigo{RIT.}
\strofa
\rigo{> Indice}
\end{multicols}

\cantotesta{76}{Alleluia la tua festa}{Canti liturgici · Alleluia}{}{}{}
\begin{multicols}{2}
\rigo{Rit.}
\rigo{A\ch{Do}{llelu}\ch{Fa}{ia} alleluia}
\rigo{A\ch{Do}{llelu}\ch{Sol}{ia} alleluia}
\rigo{A\ch{Do}{llelu}\ch{Fa}{ia} alleluia al\ch{Do}{lel}\ch{Sol}{uia} (bis\ch{Fa}{)}  \ch{Sol}{}}
\rigo{L\ch{Do}{a} tua fes\ch{Fa}{ta} non deve f\ch{Sol}{inire,}\ch{Do}{}}
\rigo{non deve f\ch{Lam}{inire} e \ch{Rem}{non} fi\ch{Sol7}{nirà} (\ch{Do}{bis)}      \ch{Do7}{}}
\strofa
\rigo{Perché\ch{Fa}{} la fest\ch{Sol}{a} siamo \ch{Do}{noi}    \ch{Lam}{}}
\rigo{che cammi\ch{Fa}{niamo} ve\ch{Sol}{rso} te          \ch{Do}{}   \ch{Do7}{}}
\rigo{perché la\ch{Fa}{} festa\ch{Sol}{} siamo \ch{Do}{noi}     \ch{Lam}{}}
\rigo{canta\ch{Re7}{ndo} insieme così…      \ch{Sol4}{}       \ch{Sol}{}}
\rigo{\ch{La}{} \ch{La}{} \ch{La}{} \ch{La}{} \ch{La}{} \ch{La}{} \ch{La}{}}
\strofa
\rigo{Rit.}
\strofa
\rigo{> Indice}
\end{multicols}

\cantotesta{77}{Alleluia, passeranno i cieli}{Canti liturgici · Alleluia}{}{}{}
\begin{multicols}{2}
\rigo{Alleluia alleluia alleluia,}
\rigo{alleluia alleluia alleluia,}
\rigo{alleluia alleluia.}
\strofa
\rigo{Passeranno i cieli e passerà la terra,}
\rigo{la sua Parola non passerà...}
\rigo{alleluia alleluia.}
\strofa
\rigo{Alleluia alleluia alleluia,}
\rigo{alleluia alleluia alleluia,}
\rigo{alleluia alleluia.}
\strofa
\rigo{> Indice}
\end{multicols}

\cantotesta{78}{Benedetto sei tu}{Canti liturgici · Offertorio}{M. Frisina}{}{}
\begin{multicols}{2}
\rigo{B\ch{La7}{ened}\ch{Re}{ett}\ch{Mim}{o} se\ch{Fa\#m}{i} Tu \ch{Sol}{Dio} \ch{Re}{del}\ch{Sim}{l'universo!}}
\rigo{D\ch{Sol}{alla}\ch{La}{} Tu\ch{Re}{a} b\ch{Sim}{ontà}}
\rigo{a\ch{Sol}{bbia}\ch{La}{mo} \ch{Re}{ric}\ch{Sim}{evuto} questo pane,}
\rigo{f\ch{Sol}{rutt}\ch{La}{o} d\ch{Re}{ell}\ch{Sim}{a} terra e del nostro lavoro}
\rigo{l\ch{Sol}{o} pr\ch{Fa\#m}{esentiamo} a Te}
\rigo{p\ch{Mim}{erch}\ch{Re}{é} diventi per noi}
\rigo{c\ch{Mim}{ibo} \ch{La}{di} vita eterna.}
\strofa
\rigo{Benedetto sei Tu Dio dell'universo!}
\rigo{Dalla Tua bontà}
\rigo{abbiamo ricevuto questo vino,}
\rigo{frutto della vite e del nostro lavoro}
\rigo{lo presentiamo a Te}
\rigo{perché diventi per noi}
\rigo{bevanda di salvezza.}
\strofa
\rigo{B\ch{Re}{ene}\ch{Mim}{dett}\ch{Fa\#m}{o} sei Tu Signor!}
\rigo{B\ch{Sol}{ened}\ch{La}{ett}\ch{Re}{o} sei Tu Signor!}
\strofa
\rigo{> Indice}
\end{multicols}

\cantotesta{79}{Benedici, o Signore}{Canti liturgici · Offertorio}{Gen Rosso}{}{}
\begin{multicols}{2}
\rigo{\ch{Rem}{Nebbia} e freddo, giorni lunghi e amari      \ch{Do}{}}
\rigo{mentre il seme muore.\ch{Rem}{}}
\rigo{P\ch{Fa}{oi} il prodigio antico e sempre nuovo   \ch{Do}{}}
\rigo{del primo filo d'erba\ch{Sib}{.}}
\rigo{E nel ven\ch{Fa}{to} dell'estate \ch{Do}{ondeggiano} \ch{Rem}{le} spighe           \ch{Fa}{}}
\rigo{avrem\ch{Do}{o} ancora pane.     \ch{La}{}       \ch{Re}{}}
\strofa
\rigo{RIT.}
\rigo{B\ch{Sol}{enedici}\ch{Re}{,} oh S\ch{Sol}{ignore,}\ch{Re}{}}
\rigo{q\ch{Do}{uesta} offerta \ch{Sol}{che} portiamo a T\ch{La}{e,}}
\rigo{f\ch{Sol}{acci} "u\ch{Re}{no"} come i\ch{Sim}{l} pane      \ch{Fa\#m}{}}
\rigo{c\ch{Mi}{he} anche oggi \ch{Sol}{hai} dato a n\ch{Re}{oi.}}
\strofa
\rigo{Nei filari dopo il lungo inverno}
\rigo{fremono le viti.}
\rigo{La rugiada avvolge nel silenzio}
\rigo{i primi tralci verdi.}
\rigo{Poi i colori dell'autunno}
\rigo{coi grappoli maturi}
\rigo{avremo ancora vino.}
\strofa
\rigo{RIT.}
\strofa
\rigo{Benedici, oh Signore,}
\rigo{questa offerta che portiamo a Te,}
\rigo{facci "uno" come il vino}
\rigo{che anche oggi hai dato a noi.}
\strofa
\rigo{> Indice}
\end{multicols}

\cantotesta{80}{Ogni mia parola}{Canti liturgici · Offertorio}{}{}{}
\begin{multicols}{2}
\rigo{Come la pioggia e la neve, scendono giù dal cielo}
\rigo{e non vi ritornano senza irrigare e far germogliare la terra,}
\rigo{così ogni mia parola non ritornerà a me,}
\rigo{senza operare quanto desidero,}
\rigo{senza aver compiuto ciò per cui l'avevo mandata.}
\rigo{Ogni mia parola. Ogni mia parola.}
\strofa
\rigo{> Indice}
\end{multicols}

\cantotesta{81}{Se m'accogli}{Canti liturgici · Offertorio}{P. Sequeri}{}{}
\begin{multicols}{2}
\rigo{\ch{Do}{Tra} le mani non ho nien\ch{Fa}{te} so che Tu mi\ch{Do}{} accogliera\ch{Mi}{i}     \ch{Lam}{}}
\rigo{chiedo so\ch{Do}{lo} di restare a\ch{Fa}{ccanto} a Te.       \ch{Sol}{}}
\rigo{Sono \ch{Do}{ricco} sola\ch{Fa}{mente} dell'amor\ch{Do}{e} che \ch{Mi}{mi} dai:    \ch{Lam}{}}
\rigo{è per que\ch{Do}{lli} che non l'hanno\ch{Fa}{} avuto mai.        \ch{Sol}{}}
\strofa
\rigo{RIT.}
\rigo{S\ch{Do}{e} m'accogli, mio Signor\ch{Mi}{e,} altro non Ti\ch{Lam}{} chiederò       \ch{Do}{}}
\rigo{e per sem\ch{Fa}{pre} la Tua stra\ch{Rem}{da} la mia strada\ch{Mi}{} resterà       \ch{Lam}{}}
\rigo{nella\ch{Fa}{} gioia e n\ch{Do}{el} dolore fino \ch{Fa}{a} quando Tu vor\ch{Mi}{rai}}
\rigo{con la ma\ch{Do}{no} nella Tu\ch{Sol}{a} camminerò.    \ch{Fa}{}      \ch{Do}{}}
\strofa
\rigo{Io Ti prego con il cuore, so che Tu mi ascolterai}
\rigo{rendi forte la mia fede più che mai.}
\rigo{Tieni accesa la mia luce fino al giorno che Tu sai,}
\rigo{con i miei fratelli incontro a Te verrò.}
\strofa
\rigo{RIT.}
\strofa
\rigo{> Indice}
\end{multicols}

\cantotesta{82}{Servo per amore}{Canti liturgici · Offertorio}{}{}{}
\begin{multicols}{2}
\rigo{\ch{Sim}{Una} notte di sudore}
\rigo{sulla barca in mezzo al mare \ch{Re}{}}
\rigo{e mentre il c\ch{La}{ielo} s'imbianca già    \ch{Fa\#m}{}}
\rigo{tu guardi\ch{Sol}{} le tue reti vuo\ch{Sim}{te.}}
\rigo{Ma la voc\ch{Re}{e} che ti chiama}
\rigo{un altro \ch{La}{mare} ti mostrerà   \ch{Re}{}}
\rigo{e sulle r\ch{Sol}{ive} di ogni cuor\ch{Re}{e}}
\rigo{le tue re\ch{Mim}{ti} g\ch{Sol}{ette}\ch{Re}{rai.}}
\strofa
\rigo{RIT.}
\rigo{O\ch{Sim}{ffri} la vita tu\ch{Re}{a}}
\rigo{come Mari\ch{La}{a} ai piedi \ch{Fa\#m}{della} croce  \ch{Sol}{}}
\rigo{e\ch{Sim}{} sarai \ch{Re}{servo} di ogni uomo,    \ch{La}{}}
\rigo{s\ch{Re}{ervo} per amore,    \ch{Sol}{}}
\rigo{s\ch{Re}{acerdote} dell'\ch{Mim}{umanità.}    \ch{Sol}{}       \ch{Re}{}}
\strofa
\rigo{Avanzavi nel silenzio}
\rigo{tra le lacrime e speravi}
\rigo{che il seme sparso davanti a te}
\rigo{cadesse sulla buona terra.}
\rigo{Ora il cuore tuo è in festa}
\rigo{perché il grano biondeggia ormai,}
\rigo{è maturato sotto il sole,}
\rigo{puoi riporlo nei granai.}
\strofa
\rigo{RIT.}
\strofa
\rigo{> Indice}
\end{multicols}

\cantotesta{83}{Beati voi}{Canti liturgici · Comunione}{}{}{}
\begin{multicols}{2}
\rigo{Beati, beati voi che mi seguite:}
\rigo{vostro sarà il regno dei cieli.}
\strofa
\rigo{Beati coloro che muoiono}
\rigo{= Alleluia beati coloro che muoiono.}
\end{multicols}

\cantotesta{84}{C'è ancora un sole}{Canti liturgici · Comunione}{}{\bEG}{}
\begin{multicols}{2}
\rigo{Ho solo un pane ma se spezzarlo tu vuoi con me}
\rigo{crescera' la letizia di marciare insieme fratel}
\rigo{Ho qui un po' d'acqua un sorso solo vuoi berlo tu ?}
\rigo{anche l'acqua di fonte a spartirlaa e' di piu'.}
\strofa
\rigo{Vecchie parole non han più senso ne voce qui,}
\rigo{sotto il fiato del vento ogni antico ricordo svanì.}
\rigo{Parole nuove sentiamo nascere in fondo al cuor,}
\rigo{sono fatte di passi di fatica, speranza e sudor.}
\strofa
\rigo{C'e' ancora un sole l'abbiamo ritrovato}
\rigo{seguiva le ombre mobili dei passi sul sentir}
\rigo{C'e' ancora un sole scaldava le tue spalle}
\rigo{quio toccai lo zaino che tu portavi per me.}
\strofa
\rigo{> Indice}
\end{multicols}

\cantotesta{85}{Canto di San Damiano}{Canti liturgici · Comunione}{}{}{}
\begin{multicols}{2}
\rigo{\ch{Re}{Ogni} u\ch{La7}{omo} sem\ch{Sim}{plice} p\ch{La7}{orta} in \ch{Re}{cuore} un sogno,\ch{La7}{}            \ch{Sim7}{}        \ch{La7}{}}
\rigo{con amore ed umiltà potrà costruirlo.}
\rigo{Se davvero tu saprai vivere umilmente}
\rigo{più felice tu sarai anche senza niente.}
\strofa
\rigo{RIT.}
\rigo{S\ch{Mi}{e} vorrai ogni \ch{Sol}{giorno} con il tu\ch{Re}{o} sudore       \ch{La7}{}}
\rigo{u\ch{Mi}{na} pietra dopo\ch{Sol}{} l'altra alto ar\ch{Re}{riverai.}       \ch{La7}{}}
\strofa
\rigo{E le gioie semplici sono le più belle,}
\rigo{sono quelle che alla fine sono le più grandi.}
\strofa
\rigo{RIT.}
\strofa
\rigo{Nella vita semplice troverai la strada}
\rigo{che la calma donerà al tuo cuore puro.}
\strofa
\rigo{RIT.}
\strofa
\rigo{> Indice}
\end{multicols}

\cantotesta{86}{Dimmelo tu}{Canti liturgici · Comunione}{}{}{}
\begin{multicols}{2}
\rigo{Il nostro c\ch{Mi}{ane} non mi riconosce p\ch{La}{iù,}\ch{Do\#m}{}     \ch{Si}{}}
\rigo{MI                              LA \#DO\#-}
\rigo{altri profumi, altre valige da portare giù,}
\rigo{c\ch{Si}{erco} le chiavi di casa,}
\rigo{m\ch{Si7}{a} questa non è casa mi\ch{Mi}{a,} più mia\ch{La}{.}}
\strofa
\rigo{I libri, quelli li ho portati via,}
\rigo{anche l'ultimo sguardo, l'ultima fotografia,}
\rigo{quella di nonna bambina}
\rigo{che gioca con il cerchio in piazza, la piazza...}
\strofa
\rigo{Amore, dimmelo tu cos\ch{Re}{'è} \ch{Fa\#m}{}     \ch{Mi}{}}
\rigo{quello che anc\ch{La}{ora} ci manca,}
\rigo{quello che ancora non c\ch{Re}{'è} \ch{Fa\#m}{}}
\rigo{e\ch{Mi}{} che ti prende alle spalle}
\rigo{e\ch{Mi7}{} non ti fa tornare indi\ch{La}{etro,} più indiet\ch{Re}{ro.}}
\rigo{T\ch{Mi}{u,} dimmelo tu cos'è\ch{La}{,}  \ch{Do\#m}{}     \ch{Si}{}}
\rigo{oh t\ch{Mi}{u,} dimmelo tu cos'è\ch{La}{.}}
\strofa
\rigo{Scopare bene, scopare bene questa è la prima cosa,}
\rigo{cercare un'altra donna un'altra casa che non sia troppo vuota}
\rigo{per ritornare di sera}
\rigo{e non sentirsi ancora più soli, ancora più soli.}
\strofa
\rigo{In questa notte d'Agosto ogni stella al suo posto lassù,}
\rigo{in questa notte di strada quella che manca sei tu,}
\rigo{in questa estate romana}
\rigo{di musica e fotografia, ancora più mia.}
\strofa
\rigo{RIT.}
\strofa
\rigo{> Indice}
\end{multicols}

\cantotesta{87}{E sono solo un uomo}{Canti liturgici · Comunione}{P. Sequeri}{}{}
\begin{multicols}{2}
\rigo{\ch{Re}{Io} lo so, S\ch{Fa\#m}{ignore} che vengo d\ch{Sol}{a} lontano,      \ch{Re}{}}
\rigo{p\ch{Re}{rima} nel pe\ch{Fa\#m}{nsiero} e poi nell\ch{Sol}{a} tua mano;      \ch{La}{}}
\rigo{i\ch{Re}{o} mi rendo co\ch{La}{nto} che Tu sei la\ch{Sol}{} mia vita          \ch{Re}{}}
\rigo{e n\ch{Sol}{on} mi sembr\ch{Sol/Fa\#}{a} vero d\ch{Mim}{i} pregarTi \ch{Mi7}{così:}      \ch{La7}{}}
\rigo{Padre d’ogni uomo e non Ti ho visto mai,}
\rigo{Spirito di vita e nacqui da una donna}
\rigo{Figlio mio fratello e sono solo un uomo}
\rigo{Eppure io capisco che Tu sei verità.}
\strofa
\rigo{Rit. E imp\ch{Re}{arerò} a guardar\ch{Sol}{e} tutto il mondo\ch{La}{}   \ch{Re7}{}}
\rigo{Sol               La                Re 7}
\rigo{con gli occhi trasparenti di un bambino}
\rigo{e insegner\ch{Sol}{ò} a chiamarTi Pa\ch{La}{dre} nostro         \ch{Re}{}       \ch{Si7}{}}
\rigo{Mi-                Mi7 (La)        La (Re)}
\strofa
\rigo{ad ogni figlio che diventa uomo. (2 v.)}
\rigo{Io lo so Signore, che Tu mi sei vicino}
\rigo{Luce alla mia mente, guida il mio cammino,}
\rigo{mano che sorregge, sguardo che perdona,}
\rigo{e non mi sembra vero che Tu esisti così.}
\rigo{Dove nasce amore Tu sei la sorgente,}
\rigo{dove c’è una Croce Tu sei la speranza,}
\rigo{dove il tempo ha fine Tu sei la vita eterna:}
\rigo{e so che posso sempre contare su di Te.}
\strofa
\rigo{Rit. E accoglierò la vita come un dono}
\rigo{E avrò il coraggio di morire anch’io}
\rigo{E incontro a Te verrò con mio fratello}
\rigo{Che non si sente amato da nessuno.}
\strofa
\rigo{> Indice}
\end{multicols}

\cantotesta{88}{Insieme}{Canti liturgici · Comunione}{}{\bEG}{}
\begin{multicols}{2}
\rigo{Insieme abbiam marciato un dì}
\rigo{per strade non battute}
\rigo{insieme abbiam raccolto un fior}
\rigo{sull'orlo di una rupe.}
\strofa
\rigo{RIT: Insieme, insieme}
\rigo{è il motto di fraternità}
\rigo{insieme, nel bene crediam.}
\strofa
\rigo{Insieme abbiam portato un dì}
\rigo{lo zaino che ci spezza}
\rigo{insieme abbiam goduto alfin}
\rigo{del vento la carezza.}
\strofa
\rigo{RIT: Insieme, insieme...}
\strofa
\rigo{Insieme abbiamo appreso ciò}
\rigo{che il libro non addita}
\rigo{abbiamo appreso che l'amor}
\rigo{è il senso della vita.}
\strofa
\rigo{RIT: Insieme, insieme...}
\strofa
\rigo{> Indice}
\end{multicols}

\cantotesta{89}{Pane del cielo}{Canti liturgici · Comunione}{}{}{}
\begin{multicols}{2}
\begin{ritornello}
\rigo{Pane del cielo}
\rigo{sei tu Gesù,}
\rigo{via d’amore:}
\rigo{tu ci fai come te.}
\strofa
\end{ritornello}
\rigo{No, non è rimasta fredda la terra:}
\rigo{tu sei rimasto con noi}
\rigo{per nutrirci di te,}
\rigo{pane di vita;}
\rigo{ed infiammare col tuo amore}
\rigo{tutta l’umanità.}
\strofa
\rigo{Sì, il cielo è qui su questa terra:}
\rigo{tu sei rimasto con noi}
\rigo{ma ci porti con te}
\rigo{nella tua casa}
\rigo{dove vivremo insieme a te}
\rigo{tutta l’eternità.}
\strofa
\rigo{No, la morte non può farci paura:}
\rigo{tu sei rimasto con noi.}
\rigo{E chi vive di te}
\rigo{vive per sempre.}
\rigo{Sei Dio con noi, sei Dio per noi,}
\rigo{Dio in mezzo a noi.}
\end{multicols}

\cantotesta{90}{Preghiera semplice}{Canti liturgici · Comunione}{}{}{}
\begin{multicols}{2}
\rigo{Se anche parlassi le lingue degli uomini e degli angeli,}
\rigo{ma non avessi la carità, sono come un bronzo che risuona,}
\rigo{o un cembalo che tintinna.}
\strofa
\rigo{E se avessi il dono della profezia,}
\rigo{e conoscessi tutti i misteri e tutta la scienza,}
\rigo{e possedessi la pienezza della Fede così da trasportare}
\rigo{le montagne, ma non avessi la carità non avrei nulla.}
\strofa
\rigo{E se anche distribuissi tutte le mie sostanze}
\rigo{e dessi il mio corpo per essere bruciato, ma non avessi la carità}
\rigo{niente mi giova.}
\strofa
\rigo{> Indice}
\end{multicols}

\cantotesta{91}{Su ali d'aquila}{Canti liturgici · Comunione}{}{}{}
\begin{multicols}{2}
\rigo{\ch{Sol7+}{Tu} che abiti al riparo\ch{Re7+}{} del Signore}
\rigo{e\ch{Sol7+}{} che dimori alla Sua ombra\ch{Re7+}{,}}
\rigo{d\ch{Fa}{i'} al Signore "Mio rifugio,\ch{Sib}{}}
\rigo{SOL7-                     LA4         LA7}
\rigo{mia roccia in cui confido".}
\strofa
\rigo{RIT.}
\rigo{E ti \ch{Re}{rialzerà,} ti so\ch{Re7+}{lleverà}}
\rigo{su ali d'\ch{Mim}{aquila,} ti r\ch{La7}{eggerà}}
\rigo{LA7-         RE7    SOL    MI-}
\rigo{sulla brezza dell'alba ti farà brillar}
\rigo{come il sole,\ch{Sim}{} cos\ch{Fa\#m}{ì} nelle S\ch{Sol}{ue} m\ch{Mim}{ani} vivrai.\ch{La7}{}            \ch{Re}{}}
\strofa
\rigo{Dal laccio del cacciatore ti libererà}
\rigo{e dalla carestia che distrugge,}
\rigo{poi ti coprirà con le Sue ali}
\rigo{e rifugio troverai.}
\strofa
\rigo{RIT.}
\strofa
\rigo{Non devi temere i terrori della notte}
\rigo{né freccia che vola di giorno,}
\rigo{mille cadranno al tuo fianco}
\rigo{ma nulla ti colpirà.}
\strofa
\rigo{RIT.}
\strofa
\rigo{Perché‚ ai Suoi angeli ha dato un comando}
\rigo{di preservarti in tutte le tue vie,}
\rigo{ti porteranno sulle loro mani}
\rigo{contro la pietra non inciamperai.}
\strofa
\rigo{RIT.}
\strofa
\rigo{E ti rialzerò, ti solleverò}
\rigo{su ali d'aquila, ti reggerò}
\rigo{sulla brezza dell'alba ti farò brillar}
\rigo{come il sole, così nelle mie mani vivrai.}
\strofa
\rigo{> Indice}
\end{multicols}

\cantotesta{92}{Symbolum 77}{Canti liturgici · Comunione}{P. Sequeri}{}{}
\begin{multicols}{2}
\rigo{\ch{Lam}{Tu} sei la mia vita, \ch{Fa}{altro} i\ch{Sol}{o} non h\ch{Do}{o;} \ch{Lam}{}}
\rigo{Tu sei la mia str\ch{Fa}{ada,} la mia\ch{Sol}{} verità,    \ch{Mi}{}}
\rigo{n\ch{Rem}{ella} Tua Parola\ch{Sol}{} io \ch{Do}{camminerò}      \ch{Lam}{}}
\rigo{f\ch{Fa}{inché‚} avrò respiro\ch{Sol}{,} fino a quando \ch{Do}{Tu} vorrai;     \ch{Mi}{}}
\rigo{n\ch{Rem}{on} avrò paura, \ch{Sol}{sai,} se Tu s\ch{Do}{ei} con me,         \ch{Lam}{}}
\rigo{i\ch{Fa}{o} Ti prego: re\ch{Sol}{sta} con me. \ch{Lam}{}}
\strofa
\rigo{Credo in Te, Signore, nato da Maria,}
\rigo{Figlio eterno e Santo, uomo come noi;}
\rigo{morto per amore, vivo in mezzo a noi,}
\rigo{una cosa sola con il Padre e con i tuoi,}
\rigo{fino a quando, io lo so, Tu ritornerai}
\rigo{per aprirci il Regno di Dio.}
\strofa
\rigo{Tu sei la mia vita, altro io non ho;}
\rigo{Tu sei la mia pace, la mia libertà,}
\rigo{niente nella vita ci separerà,}
\rigo{so che la Tua Mano forte non mi lascerà,}
\rigo{so che da ogni male Tu mi libererai}
\rigo{e nel Tuo perdono vivrò.}
\strofa
\rigo{Padre della vita, noi crediamo in Te;}
\rigo{Figlio Salvatore, noi speriamo in Te,}
\rigo{Spirito d'Amore vieni in mezzo a noi,}
\rigo{Tu da mille strade ci raduni in unità,}
\rigo{e per mille strade poi dove Tu vorrai}
\rigo{noi saremo il seme di Dio.}
\strofa
\rigo{> Indice}
\end{multicols}

\cantotesta{93}{Andate per le strade}{Canti liturgici · Finale}{}{}{}
\begin{multicols}{2}
\rigo{\ch{Sim}{Andate} per \ch{Re}{le} str\ch{La7}{ade} in t\ch{Re}{utto} il mondo}
\rigo{c\ch{Sim}{hiam}\ch{Fa\#}{ate} \ch{Sim}{i} mi\ch{La}{ei} \ch{Re}{amici} per far festa}
\rigo{c'è u\ch{Sim}{n} posto\ch{Sol}{} per\ch{La}{} ciascuno a\ch{Sim}{lla} mia mensa.}
\strofa
\rigo{Nel v\ch{Sol}{ostro} c\ch{La}{amm}\ch{Re}{ino} annunciate il Vangelo,}
\rigo{d\ch{Sim}{icen}\ch{Mim}{do:} \ch{Fa\#}{"E'} \ch{Sim}{vicino} il Regno dei cieli".}
\rigo{G\ch{Re7}{uarite} \ch{Sol}{i} malat\ch{La}{i,} mon\ch{Re}{date} i lebbrosi,}
\rigo{r\ch{Sim}{ende}\ch{Fa\#m}{te} la\ch{Do\#7}{} vita\ch{Fa\#}{} a chi l'ha perduta.}
\strofa
\rigo{Vi è stato donato con amore gratuito}
\rigo{ugualmente donate con gioia per amore.}
\rigo{Con voi non prendete nè oro nè argento,}
\rigo{perché l'operaio ha diritto al suo cibo.}
\strofa
\rigo{Entrando in una casa donatele la pace.}
\rigo{Se c'è chi vi rifiuta e non accoglie il dono,}
\rigo{la pace torni a voi ed uscite dalla casa}
\rigo{scuotendo la polvere dai vostri calzari.}
\strofa
\rigo{Ecco io vi mando come agnelli in mezzo ai lupi,}
\rigo{siate dunque avveduti come sono i serpenti,}
\rigo{ma liberi e chiari come le colombe:}
\rigo{dovrete sopportare prigioni e tribunali.}
\strofa
\rigo{Nessuno è più grande del proprio maestro}
\rigo{né il servo è più importante del suo padrone.}
\rigo{Se hanno odiato me, odieranno anche voi}
\rigo{ma voi non temete io non vi lascio soli.}
\strofa
\rigo{> Indice}
\end{multicols}

\cantotesta{94}{Con te faremo cose grandi}{Canti liturgici · Finale}{}{}{}
\begin{multicols}{2}
\rigo{Con te faremo cose grandi}
\rigo{il cammino che percorreremo insieme}
\rigo{di te riempiranno sguardi}
\rigo{la speranza che risplenderà nei volti}
\rigo{tu la voce che rischiara}
\rigo{tu la voce che ci chiama}
\rigo{tu la gioia che da vita ai nostri sogni}
\rigo{Parlaci Signore come sai}
\rigo{sei presente nel mistero in mezzo a noi}
\rigo{chiamaci col nome che vorrai}
\rigo{e sia fatto il tuo segno su di noi}
\rigo{tu la voce che rischiara}
\rigo{tu la voce che ci chiama}
\rigo{tu la gioia che da vita ai nostri sogni.}
\strofa
\rigo{Con te faremo cose grandi}
\rigo{il cammino che percorreremo insieme}
\rigo{di te riempiranno sguardi}
\rigo{la speranza che risplenderà nei volti}
\rigo{tu l'amore che dà vita}
\rigo{tu il sorriso che ci allieta}
\rigo{tu la forza che raduna i nostri giorni.}
\strofa
\rigo{Guidaci Signore come sai}
\rigo{da chi soffre chi è più piccolo di noi}
\rigo{strumenti di quel regno che tu fai}
\rigo{di quel regno che ora vive in mezzo a noi}
\rigo{tu l'amore che dà vita}
\rigo{tu il sorriso che ci allieta}
\rigo{tu la forza che raduna i nostri giorni.}
\strofa
\rigo{Con te faremo cose grandi}
\rigo{il cammino che percorreremo insieme}
\rigo{di te si riempiranno sguardi}
\rigo{la speranza che risplenderà nei volti.}
\strofa
\rigo{> Indice}
\end{multicols}

\cantotesta{95}{Fratello sole e sorella luna}{Canti liturgici · Finale}{}{}{}
\begin{multicols}{2}
\rigo{\ch{Re}{Dolce} sent\ch{Fa\#m}{ire} come\ch{Sol}{} nel\ch{La7}{} mio cuore \ch{Fa\#m}{}}
\rigo{o\ch{Sim}{ra,} \ch{La}{umi}\ch{Sol}{lmen}\ch{Fa\#m}{te,} s\ch{Sim}{ta} n\ch{Mim}{ascendo} amo\ch{Mi7}{re.}         \ch{La7}{}}
\rigo{D\ch{Re}{olce} capir\ch{Fa\#m}{e} che no\ch{Sol}{n} son pi\ch{La7}{ù} solo     \ch{Fa\#m}{}}
\rigo{m\ch{Sim}{a} che s\ch{La}{on} par\ch{Sol}{te} d\ch{Fa\#m}{i} un'\ch{Sim}{imme}\ch{Mim}{nsa} vita,   \ch{La}{}      \ch{Re}{}}
\rigo{c\ch{Sim}{he} gene\ch{Solm}{rosa} \ch{Re}{ris}\ch{Sol}{plende} \ch{Solm}{intorno} \ch{La}{a} me:      \ch{Re}{}}
\rigo{d\ch{Sim}{ono} \ch{Solm}{di} Lu\ch{Re}{i,} \ch{Sol}{del} \ch{Solm}{Suo} immenso \ch{La}{amor.}       \ch{Re}{}}
\strofa
\rigo{Ci ha dato il cielo e le chiare stelle,}
\rigo{fratello Sole e sorella Luna;}
\rigo{la madre Terra con frutti, prati e fiori}
\rigo{il fuoco, il vento, l'aria e l'acqua pura,}
\rigo{fonte di vita per le Sue creature:}
\rigo{dono di Lui, del Suo immenso amor,}
\rigo{dono di Lui, del Suo immenso amor.}
\strofa
\rigo{Sia laudato nostro Signore}
\rigo{che ha creato l'universo intero.}
\rigo{Sia laudato nostro Signore}
\rigo{noi tutti siamo Sue creature:}
\rigo{dono di Lui, del Suo immenso amor,}
\rigo{beato chi Lo serve in umiltà.}
\strofa
\rigo{> Indice}
\end{multicols}

\cantotesta{96}{Laudato sii, Signore mio}{Canti liturgici · Finale}{}{}{}
\begin{multicols}{2}
\rigo{RIT.}
\rigo{L\ch{Sol}{audato} sii, Signore mio \ch{Lam}{laudato} \ch{Re}{sii,} Signore mio.      \ch{Mim}{}}
\rigo{L\ch{Do}{audato} sii, Signore mio lau\ch{Re}{dato} si\ch{Lam}{i,} Signore \ch{Re}{mio.}           \ch{Sol}{}}
\strofa
\rigo{\ch{Sol}{Per} il sole d'ogni g\ch{Lam}{iorno}}
\rigo{c\ch{Re}{he} riscalda e dona vita,   \ch{Mim}{}}
\rigo{e\ch{Do}{gli} illumina il cammino\ch{Re}{}}
\rigo{d\ch{Lam}{i} chi cerca Te, Signore.\ch{Re}{}}
\strofa
\rigo{RIT.}
\strofa
\rigo{Per la luna e per le stelle}
\rigo{io le sento mie sorelle,}
\rigo{le hai formate su nel cielo}
\rigo{e le doni a chi è nel buio.}
\strofa
\rigo{RIT.}
\strofa
\rigo{Per la nostra madre terra}
\rigo{che ci dona fiori ed erba,}
\rigo{su di lei noi fatichiamo}
\rigo{per il pane d'ogni giorno.}
\strofa
\rigo{RIT.}
\strofa
\rigo{Per chi soffre con coraggio}
\rigo{e perdona nel Tuo amore}
\rigo{Tu gli dai la pace Tua}
\rigo{alla sera della vita.}
\strofa
\rigo{RIT.}
\strofa
\rigo{Per la morte che è di tutti}
\rigo{io la sento ogni istante,}
\rigo{ma se vivo nel Tuo amore}
\rigo{dona un senso alla mia vita.}
\strofa
\rigo{RIT.}
\strofa
\rigo{Per l'amore che è nel mondo}
\rigo{tra una donna e l'uomo suo,}
\rigo{per la vita dei bambini}
\rigo{che il mio mondo fanno nuovo.}
\strofa
\rigo{RIT.}
\strofa
\rigo{Io Ti canto, mio Signore}
\rigo{e con me la creazione}
\rigo{Ti ringrazia umilmente}
\rigo{perché Tu sei il Signore.}
\strofa
\rigo{RIT.}
\strofa
\rigo{> Indice}
\end{multicols}

\cantotesta{97}{Mani}{Canti liturgici · Finale}{}{}{}
\begin{multicols}{2}
\rigo{\ch{Re}{Vorrei} che le parol\ch{Sol}{e} mutassero \ch{La}{in} preghiera       \ch{Re}{}}
\rigo{e riv\ch{Sol}{ederti} \ch{Sim}{o} Padre che dipi\ch{Sol}{ngevi} il cie\ch{La}{lo}}
\rigo{Sapes\ch{Re}{si} quante volte\ch{Sol}{} guardando q\ch{La}{uesto} mondo        \ch{Re}{}}
\rigo{v\ch{Sol}{orrei} che tu tornass\ch{Sim}{i} a rito\ch{Sol}{ccar}\ch{La}{e} il cuore\ch{Re}{.}}
\rigo{V\ch{Sim}{orrei} che le mie mani av\ch{Fa\#m}{essero} la\ch{Sol}{} forza     \ch{Re}{}}
\rigo{per s\ch{Sol}{ostenere} ch\ch{Mi}{i} non può camminare    \ch{La}{}}
\rigo{V\ch{Sim}{orrei} che questo cuore c\ch{Fa\#m}{he} esplode in sen\ch{Sol}{timenti}         \ch{Re}{}}
\rigo{d\ch{Sol}{iventasse} c\ch{Sim}{ulla} per chi\ch{Sol}{} non ha più madr\ch{La}{e...}}
\strofa
\rigo{RIT.}
\rigo{M\ch{Re}{ani,} prendi queste mie mani,       \ch{La}{}}
\rigo{fanne vita, fanne amore  \ch{Sol}{}}
\rigo{braccia aperte per riceve\ch{Sim}{re...Chi}\ch{La}{} è solo.}
\strofa
\rigo{\ch{Re}{Cuore,} prendi questo mio cuore,    \ch{La}{}}
\rigo{fa che si spalanchi al mondo \ch{Sol}{}}
\rigo{germogliando per quegli occhi    \ch{Sim}{}}
\rigo{LA                 SOL     (LA)}
\rigo{che non sanno pianger più.}
\strofa
\rigo{Sei tu lo spazio che desidero da sempre,}
\rigo{so che mi stringerai e mi terrai la mano.}
\rigo{Fa che le mie strade si perdano nel buio}
\rigo{ed io cammini dove cammineresti Tu.}
\rigo{Tu soffio della vita prendi la mia giovinezza}
\rigo{con le contraddizioni e le falsità}
\rigo{Strumento fa che sia per annunciare il Regno}
\rigo{a chi per queste vie Tu chiami Beati...}
\strofa
\rigo{RIT.}
\strofa
\rigo{Noi giovani di un mondo che cancella i sentimenti}
\rigo{e inscatola le forze nell'asfalto di città}
\rigo{Siamo stanchi di guardare siamo stanchi di gridare}
\rigo{ci hai chiamati siamo Tuoi cammineremo insieme...}
\strofa
\rigo{RIT.}
\rigo{M\ch{Mi}{ani,} prendi queste nostre mani         \ch{Si}{}}
\rigo{fanne vita fanne amore,  \ch{La}{}}
\rigo{braccia aperte per riceve\ch{Do\#m}{re...}\ch{Si}{Chi} è solo.}
\rigo{C\ch{Mi}{uori,} prendi questi nostri cuori   \ch{Si}{}}
\rigo{fa che siano testimon\ch{La}{i}}
\rigo{che Tu chiami ogni uomo a\ch{Do\#m}{} far festa co\ch{Si}{n} Dio.      \ch{La}{}}
\strofa
\rigo{> Indice}
\end{multicols}

\cantotesta{98}{Popoli tutti}{Canti liturgici · Finale}{}{}{}
\begin{multicols}{2}
\rigo{Popoli tutti lodate il Signore, amen,}
\rigo{cantiamo amici, gloria al Signore, amen.}
\end{multicols}

\cantotesta{99}{Come Maria}{Canti liturgici · Mariani}{}{}{}
\begin{multicols}{2}
\rigo{Vogliamo vivere, Signore,}
\rigo{offrendo a Te la nostra vita;}
\rigo{con questo pane e questo vino accetta quello che noi siamo.}
\rigo{Vogliamo vivere, Signore,}
\rigo{abbandonati alla tua voce,}
\rigo{staccati dalle cose vano,}
\rigo{fissati nella vita vera.}
\strofa
\rigo{Vogliamo vivere come Maria,}
\rigo{l'irragiungibile,}
\rigo{la Madre amata,}
\rigo{che vince il mondo con l'Amore}
\rigo{e offrire sempre la Tua vita}
\rigo{che viene dal Cielo.}
\strofa
\rigo{Accetta dalle nostre mani}
\rigo{come un offerta a Te gradita,}
\rigo{i desideri di ogni cuore}
\rigo{le ansie della nostra vita.}
\rigo{Vogliamo vivere, Signore,}
\rigo{accesi dalle tue parole,}
\rigo{per ricordare ad ogni uomo}
\rigo{la fiamma viva del tuo Amore.}
\strofa
\rigo{Vogliamo vivere come Maria,}
\rigo{l'irragiungibile,}
\rigo{la madre amata}
\rigo{che vince il mondo con l'Amore}
\rigo{e offrire sempre la Tua vita}
\rigo{che viene dal Cielo.}
\rigo{e offrire sempre la Tua vita}
\rigo{che viene da Cielo.}
\strofa
\rigo{> Indice}
\end{multicols}

\cantotesta{100}{Madonna degli scout}{Canti liturgici · Mariani}{A. Ghetti - V. Ghetti}{}{}
\begin{multicols}{2}
\rigo{Madonna degli scouts, ascolta! T'invochiam:}
\rigo{concedi un forte cuore a noi, che ora partiam.}
\rigo{La strada è tanto lunga, il freddo già ci assal:}
\rigo{respingi tu, Regina, lo spirito del mal.}
\strofa
\rigo{E il ritmo dei passi ci accompagnerà:}
\rigo{là, verso gli orizzonti lontani si va.}
\strofa
\rigo{E lungo quella strada non ci lasciare tu:}
\rigo{nel volto di chi soffre saprem trovar Gesù.}
\rigo{Allor ci fermeremo le piaghe a medicar}
\rigo{e il pianto di chi è solo sapremo consolar.}
\strofa
\rigo{E il ritmo dei passi ci accompagnerà...}
\strofa
\rigo{> Indice}
\end{multicols}

\cantotesta{101}{Mamma del cielo}{Canti liturgici · Mariani}{}{\bLC}{}
\begin{multicols}{2}
\rigo{Mamma del Cielo, Vergine Santa}
\rigo{a te d’intorno il branco canta}
\rigo{Fa’ che ci amiamo come fratelli}
\rigo{mentre la pista noi percorriam.}
\strofa
\rigo{Mamma del Cielo, rendici tu}
\rigo{cortesi e buoni come Gesù.}
\strofa
\rigo{Fa che io ami ogni vicino}
\rigo{come se fosse Gesù bambino.}
\rigo{Del mio meglio fare prometto}
\rigo{per diventare un Tuo lupetto}
\strofa
\rigo{Al cuor divino stringimi tu}
\rigo{come stringevi a te Gesù.}
\strofa
\rigo{> Indice}
\end{multicols}

\cantotesta{102}{Preghiera del lupetto}{Canti liturgici · Preghiere scout}{}{\bLC}{}
\begin{multicols}{2}
\rigo{Buon Gesù}
\rigo{io ti amo con tutto il mio cuore,}
\rigo{perché mai troverò migliore amico di Te.}
\rigo{Per farti piacere io voglio,}
\rigo{con l’aiuto di Maria nostra mamma del cielo,}
\rigo{fare del mio meglio,}
\rigo{come Tu facevi ogni giorno a Nazareth.}
\rigo{Insegnami dunque ad ascoltare non me stesso}
\rigo{ma la voce,}
\rigo{quella dei miei genitori,}
\rigo{dei miei superiori,}
\rigo{dei miei capi. Così sia.}
\strofa
\rigo{> Indice}
\end{multicols}

\cantotesta{103}{Preghiera della sera}{Canti liturgici · Preghiere scout}{}{\bLC}{}
\begin{multicols}{2}
\rigo{Al cader della giornata}
\rigo{noi leviamo i cuori a Te,}
\rigo{Tu l'avevi a noi donata}
\rigo{bene spesa fu per Te.}
\strofa
\rigo{Te nel bosco, nel ruscello}
\rigo{Te nel monte Te nel mar,}
\rigo{Te nel cuore del fratello,}
\rigo{Te nel mio cercai d'amar.}
\strofa
\rigo{La strofa segreta}
\strofa
\rigo{"...e se un bel dì l'ASCI risorgerà tutti compatti ci troverà!" La guerra era finita da pochi mesi e il sogno della canzone delle Aquile randagie si avverava. Man mano che le città venivano liberate, immediatamente si ricostituivano ad opera dei vecchi Capi le prime unità e tornavano a sventolare fiamme e guidoni della risorta ASCI.}
\rigo{Così, per il S. Giorgio del 1946 il commissariato regionale lombardo fu in grado di organizzare ad Affori un grandioso raduno che attirò in un'atmosfera di entusiasmo e di fraternità molti Gruppi di altre regioni e persino dalla Svizzera. Fraterno amico, da tanti anni, dell'indimenticabile Baden e di Arnaldo Basini, fui invitato anch'io col Riparto che avevo fondato a Parma nel '45; e sebbene le condizioni dei trasporti fossero ancora molto precarie e difficili, accettai con gioia. Lo stile scout esigeva (non so se l'usanza sopravviva ancora) che chi si reca in visita a un campo porti sempre con sé qualcosa di utile da offrire: sia pure un semplice ramo per il fuoco.}
\strofa
\rigo{Così pensai di portare al campo dei fratelli milanesi, insieme con le mie quattro squadriglie, qualcosa che durasse più di un po' di legna: portai in dono le parole che avevo scritto poche settimane prima, da cantare sulla melodia della nota canzone friulana: "Ai preat la biele stele". In una notte di veglia forzata, infatti, mi erano venute spontanee e rapide, una dopo l'altra, le tre strofe di quella "Preghiera della sera" che quasi tutti gli scout, almeno parzialmente, cantano al fuoco di bivacco. Dico 'almeno parzialmente' poiché quasi tutte le raccolte di canti scout oggi non riportano la seconda, che è la più bella delle tre strofe di cui è composta la canzone "Al cader della giornata", che Baden gradì moltissimo e che volle inserita subito nel canzoniere "Canti di mezzanotte". La volete riportare sul vostro quaderno di caccia questa povera strofa derelitta?... Eccovela:}
\strofa
\rigo{Se non sempre la mia mente in Te pura si affissò;}
\rigo{se talora stoltamente da Te lungi si attardò...}
\strofa
\rigo{io Signor, ne son dolente: Te ne chiedo, o Dio, mercé!}
\rigo{el mio meglio lietamente, io farò doman per Te.}
\strofa
\rigo{Poi vengono i cieli che sembran prati e il bivacco delle stelle; ma adesso capite, se non la conoscevate ancora, perché questa strofa, a me che la scrissi per voi, sembri più bella delle altre. Voi che ne pensate?}
\strofa
\rigo{d. Tar (Aquila Azzurra)}
\rigo{da «scout-proposta educativa»}
\strofa
\rigo{Se non sempre la mia mente}
\rigo{in Te pura s'affissò;}
\rigo{se talora stoltamente}
\rigo{da Te lungi s'attardò,}
\rigo{mio Signore ne son dolente}
\rigo{Te ne chieggo, o Dio, mercè!}
\rigo{Del mio meglio lietamente}
\rigo{io farò doman per Te.}
\strofa
\rigo{I tuoi cieli sembran prati}
\rigo{e le stelle tanti fior...}
\strofa
\rigo{Son bivacchi dei beati}
\rigo{stretti in cerchio al lor Signor.}
\strofa
\rigo{Quante stelle quante stelle}
\rigo{dimmi Tu la mia qual'è'?}
\rigo{Non ambisco la più bella}
\rigo{basta sia vicino a Te!..}
\strofa
\rigo{> Indice}
\end{multicols}

\cantotesta{104}{Grazie}{Canti liturgici · Preghiere a tavola}{tradizionale}{}{}
\begin{multicols}{2}
\rigo{Grazie, signore per il tuo amore testimoniato da una croce.}
\rigo{Grazie per il sole che ci riscalda e ci fa vivere.}
\rigo{Grazie per l’amore di papà e mamma che richiama il tuo amore.}
\rigo{Grazie per la maestà delle montagne, che richiamano la tua grandezza.}
\rigo{Grazie per i piccoli fiori del prato che mi richiamano la tua fedeltà}
\rigo{Grazie per questo passero che vola nel cielo che mi richiamala libertà}
\rigo{Grazie per lo Spirito che mi illumina ogni giorno.}
\rigo{Grazie per gli amici che camminano con me.}
\rigo{Grazie per la Chiesa, madre di tutti i viventi.}
\rigo{Grazie per ogni uomo di buona volontà che cammina nelle tue vie.}
\strofa
\rigo{> Indice}
\end{multicols}

\cantotesta{105}{Per questo pan}{Canti liturgici · Preghiere a tavola}{tradizionale}{}{}
\begin{multicols}{2}
\rigo{Per questo pan, per ogni don,}
\rigo{noi ti lodiam Signor}
\rigo{Pour ce repas, pour toute joie}
\rigo{nous Te louons, Seigneur!}
\rigo{Für Hilf und Kraft, für taglich Brot}
\rigo{wir danken Dir, mein Gott!}
\rigo{For bread strength, for ev'ry gift}
\rigo{we sing praise you Lord!}
\strofa
\rigo{> Indice}
\end{multicols}

\cantotesta{106}{Al chiaror del mattin}{Canti del risveglio}{}{\bEG}{}
\begin{multicols}{2}
\rigo{Al c\ch{Sol}{hiaror} del mattin, al chiaror del mattin}
\rigo{ci sv\ch{Re7}{eglia} una canzon, ci sveglia\ch{Do}{} una canzon.   \ch{Sol}{}}
\rigo{Al chiaror del mattin, al chiaror del mattin}
\rigo{ci sv\ch{Re7}{eglia} una canzon, è sorto il sol.   \ch{Sol}{}}
\strofa
\rigo{RIT.}
\rigo{Su leviam, su leviam,}
\rigo{c'è un magnifico \ch{Re7}{ciel}\ch{Do}{o} limpido;\ch{Sol}{}}
\rigo{Su leviam, su leviam,}
\rigo{non è tem\ch{Re7}{po} di esitar\ch{Do}{.}     \ch{Sol}{}}
\strofa
\rigo{Se la pioggia cadrà, se la pioggia cadrà}
\rigo{e tutto bagnerà, e tutto bagnerà.}
\rigo{Se la pioggia cadrà, se la pioggia cadrà}
\rigo{più divertente ancor sarà.}
\strofa
\rigo{RIT.}
\strofa
\rigo{Se il vento verra', se il vento verrà}
\rigo{e tutto muovera' e tutto bagnerà.}
\rigo{Se il vento verra', se il vento verrà}
\rigo{piu' divertente ancor sara'.}
\strofa
\rigo{RIT.}
\strofa
\rigo{Se la neve verra', se la neve verrà}
\rigo{e tutto coprira', e tutto coprirà.}
\rigo{Se la neve verra', se la neve verrà}
\rigo{piu' divertente ancor sara'.}
\strofa
\rigo{RIT.}
\strofa
\rigo{> Indice}
\end{multicols}

% ---- Indici
\titolosezione{Indice alfabetico}
\begingroup\small
\begin{multicols}{2}
\lettera{A}
\voce{7}{A caccia andiamo}
\voce{68}{Acqua siamo noi}
\voce{2}{Addio lupetto}
\voce{49}{Ah io vorrei tornare}
\voce{106}{Al chiaror del mattin}
\voce{29}{Alla fiera dell'est}
\voce{76}{Alleluia la tua festa}
\voce{77}{Alleluia, passeranno i cieli}
\voce{93}{Andate per le strade}
\voce{8}{Attorno alla rupe}
\lettera{B}
\voce{83}{Beati voi}
\voce{78}{Benedetto sei tu}
\voce{79}{Benedici, o Signore}
\voce{9}{Branco di Akela}
\voce{10}{Buona caccia}
\lettera{C}
\voce{84}{C'è ancora un sole}
\voce{11}{Caccia al sambhur}
\voce{69}{Camminerò}
\voce{70}{Cantiamo a te}
\voce{50}{Canto del ricordo}
\voce{12}{Canto del risveglio}
\voce{4}{Canto della promessa}
\voce{3}{Canto della promessa del lupetto}
\voce{13}{Canto della sera}
\voce{5}{Canto delle Scolte d'Assisi}
\voce{85}{Canto di San Damiano}
\voce{39}{Canzone dell'amicizia}
\voce{40}{Caramba}
\voce{51}{Carnet di marcia}
\voce{71}{Che gioia ci hai dato}
\voce{52}{Colore del sole}
\voce{99}{Come Maria}
\voce{94}{Con te faremo cose grandi}
\lettera{D}
\voce{14}{Danza dei cani rossi}
\voce{15}{Danza del fiore rosso}
\voce{16}{Danza della tregua dell'acqua}
\voce{86}{Dimmelo tu}
\lettera{E}
\voce{87}{E sono solo un uomo}
\voce{72}{Emmanuel}
\voce{63}{Era una notte che pioveva}
\lettera{F}
\voce{53}{Faria faria oh}
\voce{17}{Fratelli alla candida luna}
\voce{95}{Fratello sole e sorella luna}
\lettera{G}
\voce{73}{Grandi cose}
\voce{104}{Grazie}
\voce{54}{Guarda lontano}
\lettera{I}
\voce{18}{Il canto del cuculo}
\voce{30}{Il duca di Barnabò}
\voce{61}{Il falco}
\voce{19}{Il lupo ferito}
\voce{20}{Il risveglio della giungla}
\voce{1}{Inno d'Italia}
\voce{88}{Insieme}
\lettera{J}
\voce{31}{Jack è in cucina}
\voce{41}{John Brown}
\lettera{K}
\voce{55}{Kamaludu}
\lettera{L}
\voce{21}{L'ultimo canto di Akela}
\voce{42}{L'uomo di Cromagnon}
\voce{32}{La battaglia di Magenta}
\voce{43}{La bella la va al fosso}
\voce{28}{La leggenda del fuoco}
\voce{33}{La macchina del capo}
\voce{56}{La mula de Parenzo}
\voce{96}{Laudato sii, Signore mio}
\voce{44}{Lazy boy}
\voce{22}{Lupo salta su}
\lettera{M}
\voce{100}{Madonna degli scout}
\voce{101}{Mamma del cielo}
\voce{97}{Mani}
\voce{45}{Me compare Giacometo}
\lettera{N}
\voce{46}{Non si va in cielo}
\lettera{O}
\voce{80}{Ogni mia parola}
\voce{47}{Oleanna}
\lettera{P}
\voce{89}{Pane del cielo}
\voce{105}{Per questo pan}
\voce{57}{Pim pam}
\voce{98}{Popoli tutti}
\voce{102}{Preghiera del lupetto}
\voce{103}{Preghiera della sera}
\voce{90}{Preghiera semplice}
\lettera{Q}
\voce{64}{Quel mazzolin di fiori}
\lettera{S}
\voce{23}{Salta o lupo}
\voce{81}{Se m'accogli}
\voce{6}{Servire}
\voce{82}{Servo per amore}
\voce{24}{Siam nella vasta giungla}
\voce{34}{Siamo andati alla caccia del leon}
\voce{62}{Signor, fra le tende schierati}
\voce{65}{Signore delle cime}
\voce{91}{Su ali d'aquila}
\voce{48}{Sul cappello un bel fior}
\voce{37}{Sul colle}
\voce{35}{Sul lago Tanganika}
\voce{66}{Sul Monte Nero}
\voce{92}{Symbolum 77}
\lettera{T}
\voce{67}{Ta-pum}
\voce{58}{Terra della betulla}
\lettera{U}
\voce{25}{Ula, ula, ula}
\voce{26}{Un bravo lupo}
\voce{36}{Un elefante}
\lettera{V}
\voce{38}{Vecchio scarpone}
\voce{59}{Vento della sera}
\voce{74}{Vivere la vita}
\voce{75}{Vocazione}
\voce{60}{Voglio girare il mondo}
\lettera{Z}
\voce{27}{Zampe agili e silenti}
\end{multicols}
\endgroup
\titolosezione{Indice per categoria}
\begingroup\small
\begin{multicols}{2}
\categoriaindice{Canti associativi}
\voce{1}{Inno d'Italia}
\categoriaindice{Canti cerimoniali}
\voce{2}{Addio lupetto}
\voce{4}{Canto della promessa}
\voce{3}{Canto della promessa del lupetto}
\voce{5}{Canto delle Scolte d'Assisi}
\voce{6}{Servire}
\categoriaindice{Canti di lupettismo}
\voce{7}{A caccia andiamo}
\voce{8}{Attorno alla rupe}
\voce{9}{Branco di Akela}
\voce{10}{Buona caccia}
\voce{11}{Caccia al sambhur}
\voce{12}{Canto del risveglio}
\voce{13}{Canto della sera}
\voce{14}{Danza dei cani rossi}
\voce{15}{Danza del fiore rosso}
\voce{16}{Danza della tregua dell'acqua}
\voce{17}{Fratelli alla candida luna}
\voce{18}{Il canto del cuculo}
\voce{19}{Il lupo ferito}
\voce{20}{Il risveglio della giungla}
\voce{21}{L'ultimo canto di Akela}
\voce{22}{Lupo salta su}
\voce{23}{Salta o lupo}
\voce{24}{Siam nella vasta giungla}
\voce{25}{Ula, ula, ula}
\voce{26}{Un bravo lupo}
\voce{27}{Zampe agili e silenti}
\categoriaindice{Fuoco di bivacco}
\sottocategoriaindice{Apertura}
\voce{28}{La leggenda del fuoco}
\sottocategoriaindice{Bans}
\voce{29}{Alla fiera dell'est}
\voce{30}{Il duca di Barnabò}
\voce{31}{Jack è in cucina}
\voce{32}{La battaglia di Magenta}
\voce{33}{La macchina del capo}
\voce{34}{Siamo andati alla caccia del leon}
\voce{35}{Sul lago Tanganika}
\voce{36}{Un elefante}
\sottocategoriaindice{Crescendo}
\voce{37}{Sul colle}
\voce{38}{Vecchio scarpone}
\sottocategoriaindice{Picco}
\voce{39}{Canzone dell'amicizia}
\voce{40}{Caramba}
\voce{41}{John Brown}
\voce{42}{L'uomo di Cromagnon}
\voce{43}{La bella la va al fosso}
\voce{44}{Lazy boy}
\voce{45}{Me compare Giacometo}
\voce{46}{Non si va in cielo}
\voce{47}{Oleanna}
\voce{48}{Sul cappello un bel fior}
\sottocategoriaindice{Calando}
\voce{49}{Ah io vorrei tornare}
\voce{50}{Canto del ricordo}
\voce{51}{Carnet di marcia}
\voce{52}{Colore del sole}
\voce{53}{Faria faria oh}
\voce{54}{Guarda lontano}
\voce{55}{Kamaludu}
\voce{56}{La mula de Parenzo}
\voce{57}{Pim pam}
\voce{58}{Terra della betulla}
\voce{59}{Vento della sera}
\voce{60}{Voglio girare il mondo}
\sottocategoriaindice{Chiusura}
\voce{61}{Il falco}
\sottocategoriaindice{Ultimi}
\voce{62}{Signor, fra le tende schierati}
\categoriaindice{Canti alpini}
\voce{63}{Era una notte che pioveva}
\voce{64}{Quel mazzolin di fiori}
\voce{65}{Signore delle cime}
\voce{66}{Sul Monte Nero}
\voce{67}{Ta-pum}
\categoriaindice{Canti liturgici}
\sottocategoriaindice{Ingresso}
\voce{68}{Acqua siamo noi}
\voce{69}{Camminerò}
\voce{70}{Cantiamo a te}
\voce{71}{Che gioia ci hai dato}
\voce{72}{Emmanuel}
\voce{73}{Grandi cose}
\voce{74}{Vivere la vita}
\voce{75}{Vocazione}
\sottocategoriaindice{Alleluia}
\voce{76}{Alleluia la tua festa}
\voce{77}{Alleluia, passeranno i cieli}
\sottocategoriaindice{Offertorio}
\voce{78}{Benedetto sei tu}
\voce{79}{Benedici, o Signore}
\voce{80}{Ogni mia parola}
\voce{81}{Se m'accogli}
\voce{82}{Servo per amore}
\sottocategoriaindice{Comunione}
\voce{83}{Beati voi}
\voce{84}{C'è ancora un sole}
\voce{85}{Canto di San Damiano}
\voce{86}{Dimmelo tu}
\voce{87}{E sono solo un uomo}
\voce{88}{Insieme}
\voce{89}{Pane del cielo}
\voce{90}{Preghiera semplice}
\voce{91}{Su ali d'aquila}
\voce{92}{Symbolum 77}
\sottocategoriaindice{Finale}
\voce{93}{Andate per le strade}
\voce{94}{Con te faremo cose grandi}
\voce{95}{Fratello sole e sorella luna}
\voce{96}{Laudato sii, Signore mio}
\voce{97}{Mani}
\voce{98}{Popoli tutti}
\sottocategoriaindice{Mariani}
\voce{99}{Come Maria}
\voce{100}{Madonna degli scout}
\voce{101}{Mamma del cielo}
\sottocategoriaindice{Preghiere scout}
\voce{102}{Preghiera del lupetto}
\voce{103}{Preghiera della sera}
\sottocategoriaindice{Preghiere a tavola}
\voce{104}{Grazie}
\voce{105}{Per questo pan}
\categoriaindice{Canti del risveglio}
\voce{106}{Al chiaror del mattin}
\end{multicols}
\endgroup

% ---- Ultima pagina: QR code verso la web app
\clearpage
\thispagestyle{empty}
\noindent
\begin{minipage}[c]{0.38\textwidth}
  \centering
  \includegraphics[width=44mm]{qr.pdf}
\end{minipage}\hfill
\begin{minipage}[c]{0.58\textwidth}
  {\fontsize{16}{19}\selectfont\bfseries Ascolta i canti e scarica gli spartiti\par}
  \vspace{3mm}
  Inquadra il codice con la fotocamera del telefono: si apre il canzoniere
  online, con l'audio di ogni canto, gli spartiti in PDF e la modalit\`a notte
  per il bivacco.\par
  \vspace{4mm}
  {\small\textbf{In caso di problemi con il QR\ldots}\par
  \begin{enumerate}\itemsep=0pt\topsep=1mm
    \item apri il browser del telefono (Chrome, Safari\ldots);
    \item scrivi nella barra degli indirizzi:\\
          {\small\url{https://lc-ncltt.github.io/canzoniere_scout/}}
    \item cerca il canto per numero (quello stampato accanto al titolo);
    \item per riaprirlo senza connessione, aggiungilo alla schermata Home.
  \end{enumerate}}
\end{minipage}

\end{document}
